% Integration fragment.  Parent supplies the common macros and theorem environments.
% This file is continuously updated; the bibliography is D_E_references.bib.
\section{D. Congruence sections and the coherent co-Lie obstruction}
\subsection{D.1. The three splitting classifications}
\status{Proved; the two exceptional nonsplitting checks have exact retained certificates.}

\begin{theorem}\label{thm:DE-split}
\status{Proved under stated hypotheses}
For $d\geq 1$ the reduction $\GL_d(\mathbb Z/p^2)\to\GL_d(\Fp)$ has a group-theoretic section precisely when
\[
 d=1\quad\text{or}\quad(d,p)=(2,2),(2,3),(3,2).
\]
For $\SL_d$ the list is
\[
 d=1\quad\text{or}\quad(d,p)=(2,3),(3,2).
\]
For $g\geq1$ the reduction $\Sp_{2g}(\mathbb Z/p^2)\to\Sp_{2g}(\Fp)$ has a section precisely when $(g,p)=(1,3)$.
\end{theorem}

\begin{proof}
The $d=1$ assertion for $\GL$ follows from the Teichm\"uller representatives; $\SL_1$ is trivial.  For $p\geq5$ let $U=1+E_{12}$.  Any lift of its reduction is $U+pX$.  In the kernel $1+p\operatorname{Mat}_d(\Fp)$, the correction to its $p$th power is the norm of $\operatorname{Ad}(U)$ applied to a matrix.  If $T=\operatorname{Ad}(U)$, then $(T-1)^3=0$, whereas
\[
 1+T+\cdots+T^{p-1}=(T-1)^{p-1}=0.
\]
Consequently $(U+pX)^p=1+pE_{12}\ne1$.  An element of order $p$ therefore cannot be lifted to an element of order $p$, excluding a section for both linear groups in every degree $d\geq2$.

Here is the completely specified finite calculation for the remaining obstructions.  Introduce a matrix variable $X_{ij}$ for each $i\ne j$ and put $\widetilde x_{ij}=1+E_{ij}+pX_{ij}$.  Impose the following necessary identities modulo $p^2$:
\[
 \widetilde x_{ij}^{\,p}=1,\qquad
 \widetilde x_{ij}\widetilde x_{kl}=\widetilde x_{kl}\widetilde x_{ij}\ (i\ne l, j\ne k),\qquad
 \widetilde x_{ij}\widetilde x_{jk}=\widetilde x_{jk}\widetilde x_{ij}\widetilde x_{ik}
 \quad(i,j,k\text{ distinct}).
\]
Because $p^2=0$, their matrix entries are affine linear equations over $\Fp$.  The retained program \texttt{code/splitting\_check.py} constructs these equations by evaluating their constant term and each single-variable increment, and performs exact modular elimination.  For $(d,p)=(3,3)$ it obtains $270$ equations in $54$ variables, coefficient rank $46$, and an inconsistent augmented system.  For $(d,p)=(4,2)$ it obtains $1920$ equations in $192$ variables, coefficient rank $177$, and an inconsistent augmented system.  These are necessary relations only: no sufficiency or presentation assumption is used for nonsplitting.  The files \texttt{certificate\_n3\_p3.json} and \texttt{certificate\_n4\_p2.json} record explicit row combinations $w$ with $wA=0$ and $wb\ne0$, and the program checks those identities.  Thus this is an exact finite proof, with a short independently checkable certificate, rather than numerical evidence.

The obstructions persist in larger degree.  Restrict a hypothetical section to the elementary matrices in a fixed $r$-dimensional block.  In every lifted matrix the off-diagonal blocks are divisible by $p$; the upper left $r\times r$ block of a product is therefore the product of the upper left blocks modulo $p^2$.  A section would consequently give matrices satisfying the contradictory $r=3,p=3$ or $r=4,p=2$ systems.  This argument applies to $\GL$ and $\SL$ alike and does not assume that the smaller blocks have determinant one.

The positive cases can also be checked exactly.  Modulo $4$, the matrices
\[
 \begin{pmatrix}0&1\\1&0\end{pmatrix},\qquad
 \begin{pmatrix}0&-1\\1&-1\end{pmatrix}
\]
generate a group of order $6$ whose reduction has order $6$.  Modulo $9$, the matrices
\[
 A=\begin{pmatrix}1&1\\6&7\end{pmatrix},\qquad
 B=\begin{pmatrix}0&4\\7&0\end{pmatrix}
\]
generate a group of order $48$ whose reduction has order $48$; $\det A=1$ and $\det B=-1$.  Its determinant-one subgroup therefore gives the section for $\SL_2(\Fp)$ at $p=3$.  For $(d,p)=(3,2)$ the six matrices, in the order $12,13,21,23,31,32$, are
\[
\begin{gathered}
\begin{pmatrix}1&1&0\\0&3&0\\0&2&3\end{pmatrix},\quad
\begin{pmatrix}1&2&1\\0&3&2\\0&0&3\end{pmatrix},\quad
\begin{pmatrix}3&0&0\\1&1&2\\0&0&3\end{pmatrix},\\
\begin{pmatrix}3&0&2\\0&3&1\\0&0&1\end{pmatrix},\quad
\begin{pmatrix}1&0&0\\2&3&0\\1&0&3\end{pmatrix},\quad
\begin{pmatrix}3&0&0\\0&3&0\\0&1&1\end{pmatrix}.
\end{gathered}
\]
They have determinant one modulo $4$ and generate a group of order $168$ whose reduction has order $168$.  The retained calculation enumerates the groups by closure under multiplication, so equality of the two orders verifies that reduction restricts to an isomorphism.

It remains to exclude $\SL_2$ at $p=2$ and the higher symplectic groups.  A determinant-one lift of $\left(\begin{smallmatrix}1&1\\0&1\end{smallmatrix}\right)$ modulo $4$ cannot have square one: writing its entries as $\left(\begin{smallmatrix}1+2a&1+2b\\2c&1+2d\end{smallmatrix}\right)$, the determinant condition gives $a+d-c=0$ modulo $2$, while the square-one condition gives $c=0$ and $a+d=1$, a contradiction.  Finally, \cite[Theorem~2.5]{DetinkoFlanneryHulpke2018} proves that for every $g\geq2$ the reduction kernel in $\Sp_{2g}(\mathbb Z/p^2)$ has no proper supplement.  A complement would be a proper supplement, so it is excluded.  Since $\Sp_2=\SL_2$, this proves the symplectic list.
\end{proof}

\begin{remark}\label{auto:D:fa384818cda6}
\status{Proved under stated hypotheses}

The distinction between a supplement and a complement is material.  The four exceptional groups in \cite[Theorem~2.5]{DetinkoFlanneryHulpke2018} are $\SL_2(\mathbb Z/4)$, $\SL_2(\mathbb Z/9)$, $\SL_3(\mathbb Z/4)$, and $\SL_4(\mathbb Z/4)$.  Their possession of proper supplements does not say that the reduction splits.  The first and fourth have no complement, as proved above.  The general-linear classification agrees with R5; the exact calculation here retains redundant relations, explaining why its equation counts differ from the earlier counts while the ranks agree.
\end{remark}

\subsection{D.2. Explicit proper examples and the failure of the universal lead}
\status{Proved: existential nonvanishing holds in every nonsplit case; universal nonvanishing is false.}

Let $\Gamma=\GL_d(\Fp)$ and let $V=\Fp^d$ be its tautological representation.  Denote by
\[
 \kappa\in\Ext^2_{\Fp[\Gamma]}(V,V)
       =H^2(\Gamma,\End_{\Fp}(V))
\]
the Postnikov class of $V\otimes^{\mathbf L}_{\mathbb Z}\Fp$, with its two cohomology modules $V$ in degrees $-1,0$.

\begin{lemma}\label{lem:DE-kappa}
\status{Proved under stated hypotheses}
The class $\kappa$ vanishes if and only if $V$ lifts to a free $\mathbb Z/p^2$-module with $\Gamma$-action.  In particular it is nonzero exactly outside the general-linear split list of Theorem~\ref{thm:DE-split}.
\end{lemma}
\begin{proof}
The change-of-rings exact sequence contains
\[
 \Ext^1_{\mathbb Z[\Gamma]}(V,V)\longrightarrow
 \Hom_{\Fp[\Gamma]}(V,V)\xrightarrow{\partial}
 \Ext^2_{\Fp[\Gamma]}(V,V),
 \qquad\partial(1_V)=\kappa.
\]
For an extension $0\to V\to L\to V\to0$, the first arrow sends its class to the map $V\to V$ induced by multiplication by $p$ on $L$.  Its value is $1_V$ precisely when $L$ is free over $\mathbb Z/p^2$ with the prescribed identification $pL\cong L/pL\cong V$.  A lift of the tautological representation is exactly a section of the general-linear reduction.
\end{proof}

Here is an explicit family of smooth projective threefolds on which this group-cohomological class survives.  The coefficients below are independent variables; thus ``generic'' specifies a field and equations, not an unspecified choice of a point of an open set.  Put $m=|\Gamma|$, $D=8m$, and
\[
 W_0=\Fp[\Gamma]^{\oplus8},\qquad
 v=(v_{h,j})_{h\in\Gamma,\,1\leq j\leq8}.
\]
The action is left translation of the index $h$.  Define homogeneous invariant polynomials $P_\alpha$ of degree $m$ by
\[
 \prod_{g\in\Gamma}\left(\sum_{h,j}z_{h,j}v_{g^{-1}h,j}\right)
    =\sum_{|\alpha|=m}z^\alpha P_\alpha(v).
\]
Take independent variables $c_{a,\alpha}$ for $1\leq a\leq D-4$ and all $\alpha$, set
\[
 k=\overline{\Fp(c_{a,\alpha})},\qquad
 Y=V\left(\sum_\alpha c_{a,\alpha}P_\alpha\ :\ 1\leq a\leq D-4\right)
      \subset\mathbb P(W_0)_k,\qquad S=Y/\Gamma.
\]

\begin{lemma}\label{lem:DE-free-ci}
\status{Proved under stated hypotheses}
$Y$ is a smooth connected projective complete-intersection threefold, the $\Gamma$-action on $Y$ is free, and $H^1(Y,\OO_Y)=0$.
\end{lemma}
\begin{proof}
For a nonidentity $g$ of order $r$, each eigenspace in $W_0\otimes\overline{\Fp}$ has dimension $D/r\leq D/2$.  This assertion remains true when $p\mid r$: a permutation cycle of length $r$ has one eigenvector for each eigenvalue which occurs, although its Jordan blocks may be nontrivial.  The nonfree locus $B$ in projective space is consequently a finite union of projective linear subspaces of dimension at most $D/2-1$.

The linear system spanned by the $P_\alpha$ is base point free.  Indeed, for a nonzero $v$, choose a linear form nonzero at every $gv$; its norm is a member of this system nonzero at $v$.  It also separates tangent directions at every point with trivial projective stabilizer.  To see this, write $L_g(z)=\sum z_{h,j}v_{g^{-1}h,j}$ and $M_g(z)=\sum z_{h,j}w_{g^{-1}h,j}$ for a tangent vector represented by $w$.  If its derivative under the projective norm map is zero, then
\[
 \sum_g M_g\prod_{h\ne g}L_h=a\prod_gL_g
\]
for a scalar $a$.  The $L_g$ are pairwise nonassociate.  Reduction modulo $L_g$ implies $L_g\mid M_g$.  In particular $w$ is proportional to $v$, so the projective tangent vector was zero.

The incidence proof of Bertini now applies on the free locus: imposing a hypersurface and its vanishing derivative gives the expected codimension because values and tangent directions are separated.  Successive independent general members meet transversally there.  Since $D-4>\dim B$, their common zero locus avoids $B$.  It is therefore a smooth complete intersection of dimension three.  The usual Koszul resolution of a positive-degree complete intersection in projective space gives $H^0(Y,\OO_Y)=k$ and $H^1(Y,\OO_Y)=0$.  The first equality proves connectedness.  No averaging by $|\Gamma|$, semisimplicity, or tame-action hypothesis has entered the argument.
\end{proof}

\begin{theorem}\label{thm:DE-coherent-obstruction}
\status{Proved under stated hypotheses}
For every pair $(d,p)$ outside the general-linear split list, the explicitly defined smooth projective threefold $S=Y/\Gamma$ above carries a rank-$d$ \etale $\Fp$-local system with full monodromy $\Gamma=\GL_d(\Fp)$ whose coherent co-Lie class $\xi$ is nonzero.  In the split cases this class vanishes on every base.  Full monodromy alone does not imply nonvanishing, even on smooth projective bases.
\end{theorem}
\begin{proof}
The quotient is a smooth projective $k$-scheme and $Y\to S$ is a connected finite \emph{\'etale} $\Gamma$-torsor.  In particular the corresponding monodromy map $\pi_1^{\et}(S)\to\Gamma$ is surjective.  Let $M=Y\times^\Gamma V$ and $E=M\otimes_{\Fp}\OO_S$.  The co-Lie complex of $D(M)$ is
\[
 \ell_{D(M)/S}\simeq M\otimes^{\mathbf L}_{\mathbb Z}\OO_S,
\]
as in \cite[Proposition~B.4]{BerghRydh2019}; its $k$-invariant is the image of $\kappa$.  The coherent Cartan--Leray spectral sequence is
\[
 H^a\bigl(\Gamma,H^b(Y,\OO_Y)\otimes_k\End_k(V\otimes k)\bigr)
       \Longrightarrow H^{a+b}(S,\End(E)).
\]
Since $H^0(Y,\OO_Y)=k$ and $H^1(Y,\OO_Y)=0$, its edge map in degree two is injective.  Extension of scalars from $\Fp$ to $k$ also preserves group cohomology and nonzero classes.  Lemma~\ref{lem:DE-kappa} therefore gives
\[
 \xi\ne0\quad\text{for every nonsplit pair }(d,p).
\]
This answers the existential question with explicit equations and a proof of survival of the class.  In the split cases the universal class is zero and hence $\xi=0$ on every base.  The last assertion follows from the curve construction below.
\end{proof}

\begin{remark}[Why Lead D is not a universal criterion]\label{rem:DE-universal-negative}
\status{Proved under stated hypotheses}
Repeat the same construction with $D-2$ independent invariant equations.  The resulting $Y$ is a smooth connected projective curve with a free $\Gamma$-action; its quotient $C$ has a local system with full monodromy $\Gamma$.  Nevertheless
\[
 H^2(C,\End(E))=0,
\]
so its coherent co-Lie class vanishes, even when $\kappa\ne0$.  Surjectivity of a fundamental-group representation is not an injectivity assertion on $H^2$, and passing from an \etale local system to its associated coherent bundle can kill the obstruction.  The Artin--Schreier sequence
\[
 0\to\End(M)\to\End(E)\xrightarrow{F-1}\End(E)\to0
\]
also identifies the kernel of the map on degree-two cohomology as the image of the connecting map from $H^1(S,\End(E))/(F-1)$; it does not force that kernel to be zero.
\end{remark}

\section{E. Further open steps}
\subsection{E.1. The two rules do determine a divided-power structure}
\status{Proved, including $p=2$.}

\begin{theorem}\label{thm:DE-PD-pi}
\status{Proved under stated hypotheses}
Let $A$ be a $\mathbb Z_{(p)}$-algebra and $I\subset A$ an ideal.  Divided-power structures on $I$ correspond bijectively, by $\gamma\mapsto\gamma_p$, to maps $\pi:I\to I$ satisfying
\[
 \pi(ax)=a^p\pi(x),\qquad
 \pi(x+y)=\pi(x)+\pi(y)+\sum_{i=1}^{p-1}\frac{x^iy^{p-i}}{i!(p-i)!}.
\]
The divided powers are given by the formula below, for every prime $p$, including $p=2$.
\end{theorem}
\begin{proof}
The apparently missing identity $p!\pi(x)=x^p$ follows from these two rules alone.  Repeated addition yields, for every positive integer $r$,
\[
 \pi(rx)=r\pi(x)+\frac{r^p-r}{p!}x^p,
\]
where the displayed coefficient is an integer.  Homogeneity gives $\pi(rx)=r^p\pi(x)$.  With $r=p+1$, therefore,
\[
 c\bigl(p!\pi(x)-x^p\bigr)=0,\qquad
 c=\frac{(p+1)^p-(p+1)}{p!}\in\mathbb Z_{(p)}^\times.
\]
Indeed the numerator has $p$-adic valuation one, also when $p=2$ (when $c=3$).  Cancellation of the unit $c$ is valid in every $A$, with arbitrary $p$-torsion.

Now \cite[Lemma~23.5.3, Tag~07GS]{StacksProject} gives existence and uniqueness of the PD structure with $\gamma_p=\pi$.  In the notation of the page, if $n=\sum_i a_ip^i$ with $0\leq a_i<p$, it is explicitly
\[
 \gamma_n(x)=\frac{(p!)^{v_p(n!)}}{n!}\prod_i\bigl(\pi^{\circ i}(x)\bigr)^{a_i},\qquad
 \gamma_0(x)=1.
\]
The coefficient is a $p$-local unit, because $v_p(n!)=\sum_i a_i(p^i-1)/(p-1)$.  This proves both necessity and sufficiency of the two rules in R11; no extra condition at $p=2$ is needed.  For completeness, successive blocks of $p$ factors and Wilson's identity give
\[
 p^{-v_p(n!)}n!\equiv(-1)^{v_p(n!)}\prod_i a_i!\pmod p,
 \qquad c_n\equiv\prod_i(a_i!)^{-1}\pmod p.
\]
Thus the page's question whether the residues are always $1$ has a negative answer for $p\geq3$ (take $n=2$), and a positive answer for $p=2$.
\end{proof}

\subsection{E.2. Crystals versus intrinsic ordinary stratifications}
\status{Refuted as stated: Theorem~\ref{thm:DE-E2-counterexample} gives a finite free relative example over a non-Noetherian base on which restriction is not faithful even for the rank-one structure crystal.  The following comparison theorems give positive replacement results; the manuscript's simultaneous locally Noetherian and coherent refinement is proved in Theorem~\ref{thm:DE-E2-coherent-Noetherian}.}

The quantifiers matter here.  The question allows arbitrary $\OO_S$-modules and a flat morphism locally of finite presentation $S\to R$.  A theorem for vector bundles over a field does not settle that question.  Conversely, failure of a particular nilpotent thickening to retract onto $S$ only defeats the usual smooth proof; it is not a counterexample to the equivalence of categories.

\begin{proposition}\label{prop:DE-E2-coherent}
\status{Proved under stated hypotheses}
Let $k$ be a field of characteristic zero and let $X$ be a scheme locally of finite type over $k$.  Restriction gives an equivalence between coherent crystals on the ordinary infinitesimal site $(X/k)_{\rm inf}$ and coherent $\OO_X$-modules with ordinary stratification on all infinitesimal neighbourhoods of the diagonal.  All the objects in these two categories are locally free of finite rank.
\end{proposition}
\begin{proof}
First, a coherent stratified module $M$ is locally free.  This can be checked after the faithfully flat extension $k\subset\bar k$.  At a closed point $x$ of an affine open $\Spec A\subset X_{\bar k}$, pull the stratification on the $n$th diagonal neighbourhood back along the second factor $x$.  It gives compatible isomorphisms
\[
 M_x/\mathfrak m_x^{n+1}M_x
  \simeq (\OO_{X,x}/\mathfrak m_x^{n+1})
                  \otimes_{\bar k}(M\otimes\kappa(x)).
\]
Consequently $\widehat M_x$ is finite free over $\widehat{\OO}_{X,x}$.  Faithful flatness of Noetherian completion implies that $M_x$ is free.  The nonfree locus of a coherent module is closed and, if nonempty on a scheme locally of finite type over $\bar k$, contains a closed point.  Thus it is empty.  The same local-freeness assertion for coherent crystals is \cite[Theorem~6.8(4), with Axiom~A of~5.2]{ToniniZhang2019}; the small and big sites give the same finite crystals, as can also be seen by the smooth-ambient construction described below.

Put $X_{\dR}(T)=X(T_{\rm red})$.  Gaitsgory--Rozenblyum's \cite[Proposition~3.4.3]{GaitsgoryRozenblyum2014} identifies derived quasi-coherent crystals with the totalization of derived quasi-coherent sheaves on the intrinsic completed \v Cech nerve
\[
             X^{\bullet+1}_{\widehat\Delta}
                      \rightrightarrows X .
\]
The hypothesis of that proposition holds for the classical scheme $X$.  An ordinary stratification on the vector bundle $M$ gives an object of this derived totalization: every pullback of a vector bundle has no higher Tor, and its ordinary isomorphisms and cocycle are also the required derived ones.  There is no unrecorded coherence obstruction, since the mapping spaces between vector bundles on an affine formal neighbourhood have no positive homotopy groups.  The cited equivalence therefore gives a derived crystal $\mathcal F$ restricting to $M$ on $X$.

We verify that $\mathcal F$ is an ordinary vector-bundle crystal; this is the essential step that prevents a derived comparison from being mistaken for the desired classical statement.  The following elementary derived nilpotent-Nakayama observation applies.  Let $B$ be a ring, let $I\subset B$ satisfy $I^N=0$, and let $K\in D(B)$ satisfy
\[
                        K\otimes_B^{\mathbf L}B/I\simeq(B/I)^r.
\]
For each $j$, associativity of derived tensor product identifies
\[
 K\otimes_B^{\mathbf L}(I^j/I^{j+1})
      \simeq (I^j/I^{j+1})^r
\]
in degree zero.  The finite filtration of $B$ by the $I^j$ then shows that $K$ itself is concentrated in degree zero.  Lift a basis modulo $I$ to obtain $B^r\to K$.  Its cokernel is zero by nilpotence.  Its kernel modulo $I$ is zero, since the assumed derived reduction also gives $\operatorname{Tor}_1^B(K,B/I)=0$.  Nilpotence again kills the kernel.  Thus $K\simeq B^r$.

For every affine infinitesimal object $U\hookrightarrow T$, derived cartesianness gives
\[
             \mathcal F_T\otimes_{\OO_T}^{\mathbf L}\OO_U
                            \simeq M|_U .
\]
Locally trivialize the vector bundle on the right and apply the observation.  It follows that $\mathcal F_T$ is an ordinary vector bundle.  Every derived transition isomorphism is therefore an ordinary pullback isomorphism, and $\mathcal F$ yields the required coherent crystal.

For completeness, finite ordinary crystals embed in the derived category just used, and restriction is fully faithful on them.  Locally embed $X$ in a smooth affine $k$-scheme $Y$.  The values of a crystal on the thickenings $X\hookrightarrow Y_n$ are vector bundles.  Any affine finite-type test scheme $T$ with a map $T_{\rm red}\to X$ admits a lift to $Y$; its defining ideal lands in a nilpotent ideal, so this lift factors through some $Y_n$.  Pull back the given value there.  Two lifts are identified by evaluating the crystal on a sufficiently large diagonal neighbourhood in $Y\times_kY$, and the triple neighbourhood gives the cocycle.  This constructs compatible vector bundles on all these classical affine tests.  By \cite[Corollary~2.2.2]{GaitsgoryRozenblyum2014} they determine a derived crystal.  The construction applies equally to morphisms and is inverse to restriction to the ordinary infinitesimal site.  Proposition~3.4.3 then proves full faithfulness and uniqueness of the extension above.  The local constructions glue, completing the proof.
\end{proof}

\par\smallskip\noindent\textbf{What this does not prove.}
The local-freeness argument uses coherence.  For example, even on the smooth affine line over a characteristic-zero field, $k[x,x^{-1}]/k[x]$ is a stratified module which is not $k[x]$-flat.  Its stratification is the quotient of the usual ones on $k[x,x^{-1}]$ and $k[x]$; infinitesimally equal functions are simultaneously invertible, so the localization has the canonical stratification.  Thus arbitrary coefficients cannot be replaced by vector bundles.  Nor does flatness of $X/R$ justify commuting arbitrary tensor products with inverse limits of principal parts.  A proof using the completed diagonal must address that issue explicitly.

Tonini--Zhang's Introduction reports an unpublished result of Bhatt for their finite-vector-bundle categories over a field.  It is not a published theorem covering the coefficients and relative bases in E.2.  Their Examples~6.14 and~6.16 provide genuine failures in a different scope: they use $k=\Fp(z)$ and $k^{\rm perf}$, which is not of finite presentation over $k$.  Those examples therefore do not refute E.2 as asked.  Finally, the derived assertion of Gaitsgory--Rozenblyum is stated in characteristic zero; neither arbitrary characteristic nor a general relative base can be silently supplied to it.


\par\smallskip\noindent\textbf{An arbitrary-coefficient positive case in characteristic $p$.}
\status{Proved and independently audited.  Coefficients are arbitrary quasi-coherent modules, and the additional hypothesis is purity of Frobenius.  This does not settle the full relative question.}

Here purity of a ring map $A\to B$ means that $N\to B\otimes_A N$ is injective for every $A$-module $N$.  It is a condition on modules, distinct from a morphism of schemes being radicial.  A scheme of characteristic $p$ is $F$-pure if its absolute Frobenius is pure on every affine open.

\begin{lemma}\label{lem:DE-crystal-Fdiv}
\status{Proved under stated hypotheses}
Let $k$ be a perfect field of characteristic $p$, and let $X$ be locally of finite type over $k$.  Quasi-coherent crystals on $(X/k)_{\rm inf}$ are equivalent to systems of arbitrary quasi-coherent modules
\[
       M_n\in\operatorname{QCoh}(X^{(n)}),\qquad
       F_{X^{(n)}/k}^*M_{n+1}\simeq M_n\quad(n\geq0).
\]
No smoothness, local freeness, coherence, or $F$-purity of $X$ is required for this equivalence.
\end{lemma}
\begin{proof}
We give the construction, including effectiveness on all infinitesimal objects.  Work first on an affine open and embed $X$ as a closed subscheme of a smooth affine $k$-scheme $Y$.  Write $F_Y^n:Y\to Y^{(n)}$ for the iterated relative Frobenius and put
\[
       Y_n=Y\times_{Y^{(n)}}X^{(n)},\qquad
       f_n:Y_n\longrightarrow X^{(n)}.
\]
In coordinates, if $X$ is defined in $Y$ by $I$, then $Y_n$ is defined by $I^{[p^n]}$.  The inclusion $X\hookrightarrow Y_n$ is nilpotent.  Since $Y$ is smooth over the perfect field, $F_Y^n$ is finite faithfully flat, and so is $f_n$.

For a quasi-coherent crystal $\mathcal E$, put $N_n=\mathcal E_{Y_n}$.  The schemes
\[
       Y_n\times_{X^{(n)}}Y_n,\qquad
       Y_n\times_{X^{(n)}}Y_n\times_{X^{(n)}}Y_n
\]
are nilpotent thickenings of the diagonally embedded $X$: they are of finite type, and relative Frobenius is a universal homeomorphism.  The crystal transition maps give an ordinary descent isomorphism for $N_n$ along $f_n$ and its cocycle condition.  Faithfully flat descent, with no finiteness restriction on the module, gives a unique quasi-coherent $M_n$ on $X^{(n)}$ with $f_n^*M_n\simeq N_n$, including its descent datum.

The inclusion $Y_n\hookrightarrow Y_{n+1}$ induces
\[
       N_n\simeq N_{n+1}|_{Y_n},\qquad
       f_{n+1}|_{Y_n}=F_{X^{(n)}/k}\circ f_n.
\]
Its maps on the double and triple fibre products show that this isomorphism respects the descent data.  Descending it along $f_n$ gives
$M_n\simeq F_{X^{(n)}/k}^*M_{n+1}$.  In particular $M_0=\mathcal E_X$.

Conversely, let $(M_n)$ be such a system and let $U\hookrightarrow T$ be an affine infinitesimal object, with ideal $J$.  For $n$ sufficiently large, every $a\in J$ satisfies $a^{p^n}=0$.  The relative Frobenius of $T$ therefore factors uniquely as
\[
       F_{T/k}^n:T\xrightarrow{\rho_n}U^{(n)}
                        \hookrightarrow T^{(n)}.
\]
Define $\mathcal E_T=\rho_n^*(M_n|_{U^{(n)}})$.  The given transition isomorphisms, together with
$\rho_{n+1}=F_{U^{(n)}/k}\circ\rho_n$, make this independent of $n$ up to a canonical compatible isomorphism.  Naturality of relative Frobenius gives all crystal pullback isomorphisms.  The construction is quasi-coherent on each $T$ and glues for nonaffine objects and open subsets.

These constructions are inverse.  In one direction, the ideal of $X$ in $Y_n$ is killed elementwise by the $p^n$th power, and the factor $\rho_n:Y_n\to X^{(n)}$ is exactly $f_n$.  Thus the construction from $(M_n)$ evaluates to $f_n^*M_n$ on $Y_n$, with its original descent datum.

In the other direction, take an arbitrary $U\hookrightarrow T$ and an $n$ as above, and restrict $Y_n$ to $U^{(n)}$.  Form
\[
       Z=Y_n\times_{X^{(n)},\rho_n}T.
\]
The map $Z\to T$ is finite faithfully flat and a universal homeomorphism.  The map $U\to Z$ induced by the two given immersions is a nilpotent closed immersion.  This last assertion can be checked without a Noetherian hypothesis on $T$: after passing to $T/J=U$, the extra ideal is a finitely generated Frobenius-kernel ideal, whose generators have $p^n$th power zero; adding the nilpotent ideal $J$ keeps it nilpotent.  Hence $Z$ is an object of the same infinitesimal site, and the crystal identifies the pullbacks of $\mathcal E_T$ and $\rho_n^*M_n$ to $Z$, compatibly with their descent data on $Z\times_T Z$.  Faithfully flat descent identifies them on $T$.  The same argument works for morphisms.  These canonical identifications also show that the construction is independent of the smooth ambient embedding and glues over $X$.
\end{proof}

\begin{proposition}\label{prop:DE-E2-Fpure}
\status{Proved under stated hypotheses}
Let $k$ be a perfect field of characteristic $p$, and let $X$ be locally of finite type over $k$ and $F$-pure.  Restriction is an equivalence between quasi-coherent crystals on $(X/k)_{\rm inf}$ and arbitrary quasi-coherent $\OO_X$-modules with ordinary stratification.
\end{proposition}
\begin{proof}
For an affine open $X=\Spec A$, the infinitesimal equivalence relations
\[
       R_n=X\times_{X^{(n)}}X
\]
are cofinal among the infinitesimal neighbourhoods of the diagonal.  Indeed, if $J$ is the finitely generated diagonal ideal in $A\otimes_k A$, then $R_n$ is defined by $J^{[p^n]}$.  For $s$ generators of $J$ one has
\[
       J^{s(p^n-1)+1}\subseteq J^{[p^n]}\subseteq J^{p^n}.
\]
The same assertion holds for the higher diagonal neighbourhoods.  An ordinary stratification is consequently exactly a compatible family of ordinary descent data for the maps $F_X^n:X\to X^{(n)}$.

By $F$-purity, these ring maps are pure, as are all their iterates.  Mesablishvili's \cite[Theorem~4]{Mesablishvili2000}, equivalently \cite[Theorem~35.4.22, Tag~08WE]{StacksProject}, proves that a pure ring map is an effective descent morphism for \emph{arbitrary} modules.  Applying it at every $n$ gives $M_n$ on $X^{(n)}$ whose pullback is the given stratified module.  Compatibility of the descent data and full faithfulness of pure descent give the isomorphisms
$F_{X^{(n)}/k}^*M_{n+1}\simeq M_n$ and all their coherences.  Thus stratified modules are equivalent to the systems in Lemma~\ref{lem:DE-crystal-Fdiv}.  That lemma completes the proof, and affine descent glues the equivalence over $X$.
\end{proof}

The purity assumption cannot simply be deleted from this proof: the same descent theorem says that a nonpure ring map is not even a descent morphism for arbitrary modules.  This observation alone is \emph{not} a counterexample to E.2, since an ineffective datum at one Frobenius level need not extend to compatible data at every higher level.  Theorem~\ref{thm:DE-E2-counterexample} supplies the missing infinite compatibility over an explicit non-Noetherian base and refutes the full assertion.  Corollary~\ref{cor:DE-E2-Noetherian-counterexample} also refutes the all-module question over a regular Noetherian base.  The simultaneous locally Noetherian and coherent-coefficient refinement is proved in Theorem~\ref{thm:DE-E2-coherent-Noetherian}, with ordinary pullbacks on both sides.


\par\smallskip\noindent\textbf{Two further positive cases, including a non-$F$-pure test.}

\begin{proposition}\label{prop:DE-E2-flat-charzero}
\status{Proved under stated hypotheses}
Let $k$ be a field of characteristic zero and let $X$ be locally of finite type over $k$.  Restriction gives an equivalence between quasi-coherent crystals whose values are flat modules and ordinary stratified quasi-coherent $\OO_X$-modules that are flat over $\OO_X$.  No finite-generation hypothesis is imposed.
\end{proposition}
\begin{proof}
A flat stratified module defines a derived object on the intrinsic completed nerve: every ordinary pullback appearing in its stratification agrees with its derived pullback and is flat over the corresponding finite diagonal neighbourhood.  The relevant mapping spaces are $0$-truncated, so the ordinary cocycle identities supply the derived coherences, exactly as for vector bundles.  Gaitsgory--Rozenblyum \cite[Proposition~3.4.3]{GaitsgoryRozenblyum2014} therefore gives a derived crystal $\mathcal F$.  We check that all its values are ordinary flat modules.

The needed nilpotent descent statement holds without any Noetherian or boundedness assumption.  Let $B$ be a ring, let $I^N=0$, and suppose $K\in D(B)$ has derived reduction $K\otimes_B^{\mathbf L}B/I=M[0]$, with $M$ flat over $B/I$.  For every $B/I$-module $Q$, associativity gives
\[
 K\otimes_B^{\mathbf L}Q\simeq M\otimes_{B/I}Q[0].
\]
Filter an arbitrary $B$-module $P$ by $I^jP$.  Its finitely many successive quotients are $B/I$-modules, so $K\otimes_B^{\mathbf L}P$ is concentrated in degree zero.  Taking $P=B$ first shows that $K$ is an ordinary module; varying $P$ then proves its flatness.  Apply this to each infinitesimal test of $X$.  Thus $\mathcal F$ is the required flat ordinary crystal.

Conversely, a flat ordinary crystal has derived Cartesian transition maps, since its values are flat.  The smooth-ambient extension to arbitrary classical affine tests used in Proposition~\ref{prop:DE-E2-coherent} applies verbatim, with flat modules in place of vector bundles.  The functorial construction also applies to morphisms.  Corollary~2.2.2 and Proposition~3.4.3 of the same source give full faithfulness and show that the two constructions are inverse.
\end{proof}

\begin{lemma}\label{lem:DE-Fdiv-universal-homeomorphism}
\status{Proved under stated hypotheses}
Over a perfect field of characteristic $p$, pullback along a finite universal homeomorphism between reduced schemes locally of finite type induces an equivalence on the categories of $F$-divided quasi-coherent modules, with arbitrary quasi-coherent coefficients.
\end{lemma}
\begin{proof}
The assertion is local on the target.  For an affine such morphism write $A\subset B$; injectivity follows because the schemes are reduced and the map is surjective.  By \cite[Lemma~29.47.9, Tag~0CNF]{StacksProject}, every $b\in B$ has a $p$-power in $A$.  Finitely many algebra generators of $B$ therefore give one $q=p^N$ with $B^q\subset A$.  The homomorphism $g^\#:B\to A$, $b\mapsto b^q$, then satisfies
\[
             f\circ g=F_X^N,\qquad g\circ f=F_Y^N.
\]
Here we use absolute Frobenius, so $g$ may be semilinear over the ground field.  Equivalently one may insert the corresponding relative Frobenius twists.  These descriptions agree because the ground field is perfect.

On $F$-divided systems, pullback by absolute Frobenius is an equivalence, with inverse the shift $(M_0,M_1,\ldots)\mapsto(M_1,M_2,\ldots)$; the structural isomorphisms give both inverse identities.  Consequently both composites $g^*f^*$ and $f^*g^*$ are equivalences, and $f^*$ is an equivalence.  Its fully faithful local inverses glue by quasi-coherent descent, proving the assertion.
\end{proof}

\begin{proposition}\label{prop:DE-E2-cusps}
\status{Proved under stated hypotheses}
Let $k$ be a perfect field of characteristic $p$, and let $X$ be an integral Cohen--Macaulay $k$-variety admitting a finite universal homeomorphism $f:Y\to X$ with $Y$ smooth.  Restriction is an equivalence between quasi-coherent crystals on $(X/k)_{\rm inf}$ and ordinary stratified quasi-coherent $\OO_X$-modules, with no finiteness or flatness requirement on the coefficients.

In particular this applies to the cusp $X=\Spec k[t^2,t^3]$, although this cusp is not $F$-pure for any prime $p$.
\end{proposition}
\begin{proof}
Ben-Zvi--Nevins \cite[Proposition~2.11, Theorem~3.11, Theorem~4.3 and Remark~4.5]{BenZviNevins2004} show that ordinary pullback $f^*$ is an equivalence on arbitrary quasi-coherent stratified modules in these hypotheses.  Their Remark~2.7 takes tensor products in the pro-category, so the stratifications used there are precisely the compatible ordinary stratifications on all finite diagonal neighbourhoods used here; no interchange with a completed tensor product is involved.  The good Cohen--Macaulay hypotheses needed for their Theorems~3.11 and~4.3 follow from smoothness of $Y$ and Proposition~2.11.

By Lemma~\ref{lem:DE-crystal-Fdiv} and Lemma~\ref{lem:DE-Fdiv-universal-homeomorphism}, $f^*$ is also an equivalence on quasi-coherent infinitesimal crystals.  The restriction functors commute with $f^*$.  The restriction functor for $Y$ is an equivalence by the smooth case, so the resulting commutative square of categories shows that the restriction functor for $X$ is an equivalence.

For the stated cusp, $k[t]$ is its finite normalization, and the map is bijective on geometric points, hence a universal homeomorphism.  The curve is a hypersurface and therefore Cohen--Macaulay.  To verify nonpurity directly, put $A=k[t^2,t^3]$ and $I=(t^2)$.  The class of $t^3$ in $A/I$ is nonzero, whereas its image under the Frobenius test map is zero:
\[
       t^3\notin(t^2),\qquad
       (t^3)^p=t^{3p}\in(t^{2p})=I^{[p]},
\]
because $t^p\in A$ for every $p\geq2$.  Thus the Frobenius map is not pure, and this positive case is not a consequence of Proposition~\ref{prop:DE-E2-Fpure}.
\end{proof}

These results by themselves do not settle the full relative E.2 assertion.  The later Theorem~\ref{thm:DE-E2-counterexample} refutes its arbitrary-module scope, while Theorem~\ref{thm:DE-E2-coherent-Noetherian} proves the simultaneous locally Noetherian and coherent refinement.  In characteristic $p$, the cusp calculation rules out that concrete non-$F$-pure family as a source of counterexamples.

\par\smallskip\noindent\textbf{A relative positive result over arbitrary bases of characteristic $p$.}

\begin{proposition}\label{prop:DE-E2-relative-Frobenius}
\status{Proved; independently audited by the parent and reduction-descent agents.}
Let $R$ be any scheme of characteristic $p$, and let $X\to R$ be locally of finite presentation.  Put $X^{(n)}=X\times_{R,F_R^n}R$.  Quasi-coherent crystals on $(X/R)_{\rm inf}$ are equivalent to systems of arbitrary quasi-coherent modules
\[
 M_n\in\operatorname{QCoh}(X^{(n)}),\qquad
 F_{X^{(n)}/R}^*M_{n+1}\simeq M_n.
\]
If every relative Frobenius $F_{X^{(n)}/R}$ is pure, restriction gives an equivalence between these crystals and ordinary stratified quasi-coherent $\OO_X$-modules.  No perfectness, reducedness, or Noetherian hypothesis on $R$ is required.
\end{proposition}
\begin{proof}
We verify that the construction of Lemma~\ref{lem:DE-crystal-Fdiv} is relative.  Locally on $X$ and $R$, choose a finitely presented closed immersion $X\hookrightarrow Y=\mathbf A_R^d$.  The iterated relative Frobenius $F_{Y/R}^n:Y\to Y^{(n)}$ is finite faithfully flat: in coordinates its module basis consists of the monomials with all exponents less than $p^n$.  More generally the same statement for smooth $Y/R$ is \cite[Lemma~50.19.5, Tag~0FW2]{StacksProject}.

Put $Y_n=Y\times_{Y^{(n)}}X^{(n)}$.  If $I=(f_1,\ldots,f_s)$ defines $X$ in $Y$, then $Y_n$ is defined by $(f_1^{p^n},\ldots,f_s^{p^n})$.  This remains true even though $F_R$ need not be flat: the ideal of a base-changed closed immersion is the image of the tensor product ideal, and relative Frobenius sends $f_i\otimes1$ to $f_i^{p^n}$.  Thus $X\hookrightarrow Y_n$ is nilpotent, and $Y_n\to X^{(n)}$ is finite faithfully flat.  Its double and triple fibre products are nilpotent thickenings of the diagonal copy of $X$.  Indeed, their defining difference ideals have finitely many generators killed by $p^n$th powers, in addition to the displayed nilpotent ideal.  A crystal therefore descends on this finite faithfully flat cover to a module $M_n$ on $X^{(n)}$.

The inclusion $Y_n\hookrightarrow Y_{n+1}$ and its maps on the double and triple products descend to $M_n\simeq F_{X^{(n)}/R}^*M_{n+1}$.  Conversely, for any nilpotent test $U\hookrightarrow T$ with ideal $J$, choose $n$ such that $J^{[p^n]}=0$.  In affine coordinates, the rule
\[
 (\OO_T/J)\otimes_{\OO_R,F_R^n}\OO_R
        \longrightarrow\OO_T,
       \qquad \overline b\otimes r\longmapsto b^{p^n}r
\]
is well defined and gives the relative Frobenius factor $\rho_n:T\to U^{(n)}$.  Set $\mathcal E_T=\rho_n^*(M_n|_{U^{(n)}})$.  The structural isomorphisms make this independent of $n$ and give the crystal transitions.

The inverse comparison is detected on
\[
                 Z=Y_n\times_{X^{(n)},\rho_n}T.
\]
It is finite faithfully flat over $T$ and is a nilpotent thickening of $U$.  For the latter assertion, modulo $J$ the extra difference ideal is finitely generated and killed elementwise by the $p^n$th power; adding $J$ preserves nilpotence.  This argument does not require $T$ to be Noetherian.  Ordinary faithfully flat module descent, including the double-product datum, proves that the constructions are inverse on all objects and morphisms.  Independence of the smooth embedding and gluing follow by the same canonical comparison.

Finally the relations $X\times_{X^{(n)}}X$ are cofinal among diagonal neighbourhoods: if $D$ is the diagonal ideal, their ideals are $D^{[p^n]}$, and the finite-generation bounds used in Proposition~\ref{prop:DE-E2-Fpure} apply over $R$ unchanged.  Under the stated purity hypotheses all iterated relative Frobenius maps are pure.  Effective module descent \cite[Theorem~35.4.22, Tag~08WE]{StacksProject} supplies $M_n$ from the stratification at every level; full faithfulness supplies the transition identifications.  This proves the second assertion.
\end{proof}

\begin{corollary}\label{cor:DE-E2-relative-node}
\status{Proved under stated hypotheses}
For every $\Fp$-algebra $R$, the flat finitely presented nodal family
\[
                    X=\Spec R[x,y]/(xy)\longrightarrow\Spec R
\]
satisfies the crystal--stratification equivalence for arbitrary quasi-coherent coefficients.
\end{corollary}
\begin{proof}
As an $R$-module its algebra is free on $1,x,x^2,\ldots,y,y^2,\ldots$, so the family is flat.  Every twist has the same displayed presentation.  The relative Frobenius $R[x_n,y_n]/(x_ny_n)\to R[x,y]/(xy)$ sends $x_n$ to $x^p$, $y_n$ to $y^p$, and fixes $R$.  It has an $R[x_n,y_n]/(x_ny_n)$-linear retraction: retain precisely the monomials whose exponent is divisible by $p$, divide that exponent by $p$, and keep their coefficients in $R$.  Linearity follows directly from $xy=0$.  Hence all relative Frobenius maps are split, in particular pure, and Proposition~\ref{prop:DE-E2-relative-Frobenius} applies.
\end{proof}

\par\smallskip\noindent\textbf{A counterexample to the full flat, locally finitely presented assertion.}

\begin{theorem}\label{thm:DE-E2-counterexample}
\status{Proved; direct all-tests argument independently audited by the parent and reduction-descent agents.}
For every prime $p$ there is an $\Fp$-algebra $R$ and a finite free morphism $X\to\Spec R$ of rank $p$ for which the restriction functor
\[
 \{\text{crystals of $\OO$-modules on $(X/R)_{\rm inf}$}\}
 \longrightarrow
 \{\text{ordinary stratified $\OO_X$-modules relative to $R$}\}
\]
is not faithful.  The failure already occurs among automorphisms of the rank-one locally free structure crystal.  Thus Part~E.2, with precisely its stated hypotheses, has a negative answer.
\end{theorem}
\begin{proof}
Put
\[
 R_0=\Fp[t],\qquad I=\Fp[t,t^{-1}]/\Fp[t],\qquad
 R=R_0\ltimes I.
\]
Here the last expression is the square-zero extension with multiplication
$(a,u)(b,v)=(ab,av+bu)$.  We regard $t$ as $(t,0)$ and $I$ as the ideal $0\oplus I$ of $R$.  For $n\geq0$ set
\[
 s_n=1+p+\cdots+p^{n-1}=\frac{p^n-1}{p-1},\qquad
 b_n=(0,[t^{-s_n}])\in I,
\]
where $s_0=0$, so that $b_0=0$.  These elements satisfy
\begin{equation}\label{eq:DE-E2-Pruefer-recurrence}
 t^{p^n}b_{n+1}=b_n\quad(n\geq0),\qquad
 tb_1=0,\qquad b_1\ne0.
\end{equation}
Let
\[
 f=x^p-tx,\qquad A=R[x]/(f),\qquad X=\Spec A.
\]
The polynomial $f$ is monic.  Hence $A$ is a free $R$-module with basis $1,x,\ldots,x^{p-1}$; in particular $X/R$ is flat and finitely presented.  Moreover $f$ is a non-zero-divisor in $R[x]$, and every fibre is a zero-dimensional hypersurface over a field, hence a complete intersection and Cohen--Macaulay.  The example does not rely on a failure of the Cohen--Macaulay property of the fibres.

We construct a nonzero endomorphism $\delta$ of the structure crystal whose evaluation on $X$ is zero.  Let $U\hookrightarrow T$ be any infinitesimal test object, with nilpotent ideal $J$.  Locally on $T$, choose an integer $n$ such that $J^{[p^n]}=0$ and a lift $z\in\OO_T$ of the function $x|_U$.  Define
\begin{equation}\label{eq:DE-E2-phantom-value}
                     \delta_T=b_nz^{p^n}\in\OO_T.
\end{equation}
Changing $z$ by an element of $J$ does not change this function, because $T$ has characteristic $p$.  Also $z^p-tz\in J$, and therefore
\[
 b_{n+1}z^{p^{n+1}}-b_nz^{p^n}
       =b_{n+1}(z^p-tz)^{p^n}=0
\]
by \eqref{eq:DE-E2-Pruefer-recurrence}.  The definition is consequently independent of every sufficiently large $n$.  These two observations show that the local functions glue, including when only local nilpotence bounds have been chosen.

For a morphism of infinitesimal test objects, choose a common sufficiently large $n$ and pull back a local lift $z$.  Formula~\eqref{eq:DE-E2-phantom-value} is then preserved by pullback.  Thus the functions $\delta_T$ define a natural $\OO$-linear endomorphism of the structure crystal.  This construction works on all test objects, without a finiteness or quasi-coherence restriction on their nilpotent ideals.  It also works for the big infinitesimal site if that convention is adopted.

On the test object $X\hookrightarrow X$ one may take $n=0$, giving
$\delta_X=b_0x=0$.  On the other hand, take the square-zero test
\[
             X\hookrightarrow T=\Spec B,\qquad B=R[x]/(f^2).
\]
Here $J=(f)/(f^2)$ satisfies $J^2=0$, so $n=1$ is allowed for every $p$.  Formula~\eqref{eq:DE-E2-phantom-value} gives
\[
                  \delta_T=b_1x^p=b_1f\ne0.
\]
The equality uses $tb_1=0$.  Nonvanishing follows from the monicity of $f^2$: the classes of $1,x,\ldots,x^{2p-1}$ form an $R$-basis of $B$, and the coefficient of $x^p$ is $b_1\ne0$.  In fact this test $T$ is itself finite free over $R$.

Finally, $I^2=0$ implies $\delta_T^2=0$ for every $T$.  Hence $1+\delta$ is an automorphism of the rank-one structure crystal, with inverse $1-\delta$, and its restriction to the ordinary stratified module $\OO_X$ is the identity.  The restriction functor is not faithful, proving the assertion.
\end{proof}


\begin{corollary}[Failure over a regular Noetherian base with arbitrary quasi-coherent coefficients]\label{cor:DE-E2-Noetherian-counterexample}
\status{Proved; independently audited by the parent and reduction-descent agents.}
The restriction functor of Theorem~\ref{thm:DE-E2-counterexample} can fail to be faithful even when the base is the regular Noetherian ring $R=\Fp[t]$ and $X/R$ is finite free with complete-intersection fibres.  In this version both crystals are quasi-coherent; one is the rank-one structure crystal and the other has arbitrary, noncoherent coefficients.
\end{corollary}
\begin{proof}
Use $R=\Fp[t]$, the $R$-module $I=R[t^{-1}]/R$, and $A=R[x]/(x^p-tx)$.  We use no square-zero extension of the base.  Besides the structure crystal, consider the constant-base coefficient crystal
\[
                 \mathcal E_T=\OO_T\otimes_R I.
\]
Its crystal transition isomorphisms are ordinary tensor associativity, so no flatness assumption on $I$ is required.  Each value is a quasi-coherent module.

With $s_n$ and $b_n=[t^{-s_n}]\in I$ as before, the sections
\[
                    z^{p^n}\otimes b_n\in\mathcal E_T
\]
define a morphism $\OO_{\rm inf}\to\mathcal E$.  The lift-independence, compatibility in $n$, gluing, and naturality on all test objects are exactly the two identities proved in Theorem~\ref{thm:DE-E2-counterexample}, now interpreted in a module rather than a square-zero ideal.  Its value on $X$ is zero, whereas its value on $B=R[x]/(x^p-tx)^2$ is
\[
                         x^p\otimes b_1\ne0.
\]
Indeed, $B\otimes_R I$ is the direct sum of $2p$ copies of $I$ in the monic basis, and its $x^p$ coefficient is $b_1\ne0$.  This proves nonfaithfulness.  Equivalently, the upper-triangular automorphism of $\OO_{\rm inf}\oplus\mathcal E$ with this off-diagonal entry is a nonidentity automorphism invisible to restriction.

The module $I$ is not finitely generated over $R$.  Since $A$ is finite faithfully flat and free over $R$, the module $A\otimes_R I$ is not finitely generated over $A$: otherwise it would be finitely generated over $R$, as would its direct summand $I$.  Thus this example does not fall under the manuscript's coherent-coefficient proviso.
\end{proof}

\begin{remark}[The exact role of the base hypothesis]\label{rem:DE-E2-Noetherian-scope}
\status{Proved under stated hypotheses}
The ring $R$ in Theorem~\ref{thm:DE-E2-counterexample} is non-Noetherian: its square-zero ideal $I$ is not finitely generated, since finitely many Laurent principal parts have a bounded pole order whereas $I$ contains every $t^{-m}$.  The supplied transcription at cote~6, scan~7 adds ``du moins si R loc.noeth., et se bornant aux Modules coh\'erents''.  The counterexample refutes the broader assertion actually requested in Part~E.2, even with finite locally free coefficients, and does not refute that inserted locally Noetherian version.  Corollary~\ref{cor:DE-E2-Noetherian-counterexample} further refutes the all-module assertion over a regular Noetherian base.  The simultaneous locally Noetherian and coherent-coefficient refinement is proved in Theorem~\ref{thm:DE-E2-coherent-Noetherian}.  This positive theorem uses both qualifications in the manuscript insertion; neither counterexample contradicts it.

There is also a direct description using Proposition~\ref{prop:DE-E2-relative-Frobenius}.  The $n$th Frobenius twist has algebra
\[
 A_n=R[x_n]/(x_n^p-t^{p^n}x_n),
\]
and multiplication by $a_n=b_nx_n$ defines an endomorphism of the unit $F$-divided system, since
\[
 F_{X^{(n)}/R}^{\#}(a_{n+1})
       =b_{n+1}x_n^p=b_nx_n=a_n.
\]
Here $a_0=0$ but $a_1\ne0$.  This is an infinitely compatible failure, not merely an ineffective datum for one Frobenius map.  The preceding direct proof is independent of the $F$-divided comparison theorem.
\end{remark}

\par\smallskip\noindent\textbf{Faithfulness in the manuscript's Noetherian coherent refinement.}

\begin{proposition}\label{prop:DE-E2-coherent-faithfulness}
\status{Proved; root and reduction-descent full independent proof/source audits passed.}
Let $S$ be locally Noetherian and let $X\to S$ be flat and locally of finite presentation.  Let $\mathcal E,\mathcal F$ be quasi-coherent crystals on the ordinary infinitesimal site of $X/S$, with $\mathcal F_X$ coherent.  Then evaluation on $X$ is injective on morphisms:
\[
 \Hom_{(X/S)_{\rm inf}}(\mathcal E,\mathcal F)
       \longrightarrow \Hom_{\OO_X}(\mathcal E_X,\mathcal F_X).
\]
In particular the restriction functor from coherent crystals to ordinary coherent stratified modules is faithful.  The assertion holds for the small and the big infinitesimal site, including when arbitrary, possibly non-Noetherian, test schemes are allowed.  No finiteness assumption on $\mathcal E$ is required.
\end{proposition}
\begin{proof}
Let $u:\mathcal E\to\mathcal F$ have $u_X=0$.  We first work over an affine base $S=\Spec R$ and an affine open $V=\Spec A\subset X$ with $A$ a finitely presented flat $R$-algebra.  Consider an affine Noetherian test $U=\Spec C\hookrightarrow T=\Spec B$, whose structural map $U\to X$ factors through $V$.  Thus $B\twoheadrightarrow C$ has a nilpotent kernel $I$, and there is an $R$-algebra map $A\to C$.  In the small site, after restricting to affine opens, we may take $U=V$.

The crystal property identifies $\mathcal F_T/I\mathcal F_T$ with the pullback of $\mathcal F_X$ to $U$, which is a finite $C$-module.  Lifting a finite set of generators and iterating $I^N=0$ shows that $\mathcal F_T$ is a finite $B$-module; this elementary nilpotent form of Nakayama's argument does not require $\mathcal F_T$ to be finite in advance.

Set
\[
 D=A\otimes_R B,\qquad
 J=\ker\bigl(D\longrightarrow C,\ a\otimes b\longmapsto a\overline b\bigr).
\]
The map to $C$ is surjective.  The ring $D$ is Noetherian, since it is a finitely presented $B$-algebra, and it is flat over $B$, since $A$ is flat over $R$.  For every $n\geq1$, the closed immersion
\[
 U\hookrightarrow T_n=\Spec(D/J^n)
\]
is an infinitesimal test, with compatible maps to $T$ and to $V$.  The latter map and $u_X=0$ give $u_{T_n}=0$.  The former map and the crystal transition isomorphisms therefore give
\[
 (D/J^n)\otimes_B u_T=0\qquad(n\geq1).
\]
Put $N=D\otimes_B\mathcal F_T$, a finite $D$-module.  Every element in the image of $D\otimes_B u_T$ belongs to $\bigcap_nJ^nN$.

Fix any prime $\mathfrak q\subset B$.  Since $I$ is nilpotent, it determines a prime of $C=B/I$ and hence a prime $\mathfrak p\subset D$ containing $J$, with $\mathfrak p\cap B=\mathfrak q$.  Krull's intersection theorem in the Noetherian local ring $D_{\mathfrak p}$ gives
\[
 \bigcap_{n\geq1}(JD_{\mathfrak p})^nN_{\mathfrak p}=0.
\]
Thus $D_{\mathfrak p}\otimes_Bu_T=0$.  The map $B_{\mathfrak q}\to D_{\mathfrak p}$ is flat and local, hence faithfully flat.  Consequently $(u_T)_{\mathfrak q}=0$.  As $\mathfrak q$ was arbitrary, $u_T=0$.  The precise commutative-algebra inputs are Krull intersection and faithful flatness of a flat local map, \citet[Lemmas~10.51.4 and~10.39.17, Tags~00IP and~00HR]{StacksProject}; the slightly more general separation statement is Lemma~10.51.5, Tag~00IQ.

It remains to remove the Noetherian condition on the test.  Locally write $A=R[z_1,\ldots,z_r]/K$, lift the images of the $z_i$ from $C$ to $B$, and choose $N$ with $I^N=0$.  The resulting map $R[z_1,\ldots,z_r]\to B$ annihilates $K^N$, and hence factors through the Noetherian ring
\[
 B_0=R[z_1,\ldots,z_r]/K^N.
\]
The test $(U,T)$ maps to the finite-type test
$V\hookrightarrow\Spec B_0$.  In the small site $U$ is an open of $V$, as required; in the big site its given map to $V$ suffices.  We already proved $u_{\Spec B_0}=0$.  The crystal transition isomorphisms then give $u_T=0$.  These local conclusions glue over $X$ and $S$, proving the proposition.
\end{proof}

\begin{remark}[What this proves in the coherent refinement]\label{rem:DE-E2-coherent-faithfulness-scope}
\status{Proved consequence of Proposition~\ref{prop:DE-E2-coherent-faithfulness}.}
The nonfaithfulness in Theorem~\ref{thm:DE-E2-counterexample} uses a non-Noetherian base, and that in Corollary~\ref{cor:DE-E2-Noetherian-counterexample} uses a noncoherent target.  The proposition shows that both cannot be avoided simultaneously.  The proposition itself asserts faithfulness, even with an arbitrary quasi-coherent source.  Fullness and essential surjectivity for coherent source and target are established separately in Theorem~\ref{thm:DE-E2-coherent-Noetherian}, using completed ordinary descent and an explicit zeroth-cohomology comparison.  A derived comparison theorem alone would not establish those ordinary conclusions.
\end{remark}

\par\smallskip\noindent\textbf{Flat coefficients over a Noetherian base, in every characteristic.}

\begin{proposition}\label{prop:DE-E2-flat-Noetherian}
\status{Proved; root and monodromy full independent proof/source audits passed.}
Let $S$ be locally Noetherian and let $X\to S$ be flat and locally of finite presentation.  Restriction is an equivalence between quasi-coherent crystals whose values on infinitesimal tests are flat modules and ordinary stratified quasi-coherent $\OO_X$-modules that are flat over $\OO_X$.  There is no finite-generation hypothesis on the coefficients, and no restriction on the characteristic of $S$.

This equivalence restricts to finite locally free coefficients.  It holds for either the small or big infinitesimal site, with arbitrary classical test schemes.
\end{proposition}
\begin{proof}
We state explicitly the relative derived comparison being used.  Zariski locally, $S$ and $X$ are affine Noetherian schemes and $X/S$ is separated, of finite type and of finite Tor-amplitude.  Jiang's Theorem~2.1 and Remark~2.4 identify the category of derived stratifications on the intrinsic formal diagonal nerve with modules over his derived differential-operator algebra.  His Theorem~6.3 identifies the latter category with derived quasi-coherent crystals on $(X/S)_{\rm inf}$, compatibly with evaluation on $X$; Appendix~B, Theorem~B.5, identifies the small and big versions \citep[pp.~16--17, 32--34 and 43--45]{Jiang2023Derived}.

For clarity, the formal-diagonal interpretation uses Jiang's Proposition~6.2: on an affine closed immersion with ideal $J$, the category of derived sheaves on the formal completion is
\[
       \varprojlim_n D(B/J^n),
\]
with derived pullback transitions.  Equivalently it is the category of derived $J$-complete objects, or, by the local-cohomology/completion equivalence, the category denoted $\Gamma_J(\mathrm{QCoh}(\Spec B))$ in his Theorem~2.1.  The pullbacks in that theorem are the corresponding supported pullbacks.  Remark~2.4 gives precisely the resulting derived stratifying category.  We use this comparison only after checking that the ordinary data define objects and morphisms in these derived categories.  Flatness of $X/S$ ensures that its iterated products over $S$ here are the ordinary products, without extra derived structure.

Let $M$ be an ordinary flat stratified $\OO_X$-module.  On every finite diagonal neighbourhood, all its pullbacks are flat over that neighbourhood.  Their ordinary pullbacks therefore equal their derived pullbacks.  The given compatible isomorphisms and cocycle conditions define an object on the derived formal diagonal nerve.  Mapping spaces between the ordinary modules at each level have no positive homotopy groups, so these ordinary cocycle identities supply the required coherent identifications.  The cited comparison supplies a derived crystal $\mathcal K$ with $\mathcal K_X=M[0]$.

We now prove that its values are ordinary flat modules; this is a separate step from the cited comparison.  Let $B$ be any ring, let $I^N=0$, and let $K\in D(B)$ satisfy
\[
          K\otimes_B^{\mathbf L}B/I\simeq P[0],
          \qquad P\text{ flat over }B/I.
\]
For every $B/I$-module $Q$, tensor associativity gives
$K\otimes_B^{\mathbf L}Q\simeq(P\otimes_{B/I}Q)[0]$.
Filter an arbitrary $B$-module $Q'$ by the finite filtration $I^jQ'$.  Each successive quotient is a $B/I$-module, so the associated exact triangles show that $K\otimes_B^{\mathbf L}Q'$ is concentrated in degree zero.  Taking $Q'=B$ proves that $K$ is an ordinary module, and varying $Q'$ proves that it is flat.  This argument allows unbounded $K$ and imposes no Noetherian condition on $B$.

For a test $U\hookrightarrow T$, apply this argument to $K=\mathcal K_T$, whose derived reduction is the flat module pulled back from $M$ on $U$.  Thus $\mathcal K$ is an ordinary flat crystal.  Conversely an ordinary flat crystal has derived Cartesian transition maps, so defines a derived crystal.  The two constructions are inverse by the cited comparison, and the same reasoning on the zero-truncated mapping spaces identifies their ordinary morphisms.

These arguments initially use Jiang's finite-type classical tests.  They extend to arbitrary tests without an assumption about infinite completed tensor products.  Indeed, locally a test $\Spec C\hookrightarrow\Spec B$ mapping to an affine $\Spec A\subset X$, with $A=R[z_1,\ldots,z_r]/K$, maps to a finite-type test
\[
 \Spec A\hookrightarrow\Spec\bigl(R[z_1,\ldots,z_r]/K^N\bigr)
\]
by lifting the $z_i$, for sufficiently large $N$.  Two such choices have a common refinement: use the finite-type product of their ambient thickenings, then a sufficiently high infinitesimal neighbourhood of the diagonal copy of $\Spec A$.  Parallel compatible maps are equalized by their closed equalizer, which still contains that copy of $\Spec A$.  Pullback of the finite-type crystal is therefore independent of all choices and gives the unique extension, on objects and morphisms.  This also proves compatibility with affine localization and gluing.

Finally, if $M$ is finite locally free, nilpotent Nakayama shows that every test value is finitely generated.  At a local test ring, lift a basis of its reduction and use flatness to identify the kernel modulo the nilpotent ideal; a second nilpotent Nakayama argument makes that kernel zero.  The test value is consequently finite locally free.  The equivalence thus restricts as asserted.
\end{proof}

\begin{corollary}\label{cor:DE-E2-coherent-any-field}
\status{Proved; root and monodromy full independent proof/source audits passed.}
For every field $k$ and every $k$-scheme $X$ locally of finite type, restriction gives an equivalence between coherent infinitesimal crystals and ordinary coherent stratified $\OO_X$-modules.  Smoothness, reducedness and $F$-purity are unnecessary.  Every coefficient module in either category is locally free of finite rank.
\end{corollary}
\begin{proof}
A coherent stratified module is locally free.  After the faithfully flat extension to an algebraic closure, fix a closed point $x$ and pull its stratification on each diagonal neighbourhood back along the second factor $x$.  The compatible isomorphisms identify the completed stalk with
\[
        \widehat{\OO}_{X,x}\otimes_k M_x/\mathfrak m_xM_x.
\]
Faithful flatness of the Noetherian local completion then makes the stalk free.  The nonfree locus of a finitely presented module is closed, so checking closed points on each affine finite-type open suffices.  Local freeness descends along the field extension.  This is also the argument behind Tonini--Zhang's Theorem~6.8(3).

For coherent infinitesimal crystals, Tonini--Zhang's Theorem~6.8(4), with its proof using Lemmas~6.10 and~6.15, shows that every test value is finite locally free \citep{ToniniZhang2019}.  The statement includes arbitrary ground fields.  The proposition now applies, since a scheme locally of finite type over a field is flat over that field, and proves the asserted equivalence.
\end{proof}

\begin{remark}[Scope and the subsequent coherent comparison]\label{rem:DE-E2-flat-Noetherian-scope}
\status{Proved scope statement.}
The proposition does not replace ordinary nonflat coherent coefficients by flat ones over a general base.  Such coefficients occur even for smooth morphisms: the pullback of a coherent module from the base has its constant-base crystal, regardless of its flatness.  Theorem~\ref{thm:DE-E2-coherent-Noetherian} below separately proves fullness and effectivity for nonflat coherent coefficients over every locally Noetherian base.  The present proposition has a different strength: it allows arbitrary, possibly infinitely generated, flat coefficients.  The original arbitrary-module assertion in Part~E.2 remains refuted by the explicit examples above.
\end{remark}

\par\smallskip\noindent\textbf{The full Noetherian coherent refinement.}

\begin{lemma}[An exact completed \v{C}ech complex over a nilpotent test]\label{lem:DE-E2-completed-Cech}
\status{Proved; root and monodromy full independent proof/source audits passed.}
Let $R$ be Noetherian, let $A$ be a flat finitely presented $R$-algebra, and let $B$ be a Noetherian $R$-algebra with a nilpotent ideal $I$.  Put $C=B/I$ and fix an $R$-algebra map $a:A\to C$.  For $q\geq0$ put
\[
 D^q=B\otimes_R A^{\otimes_R(q+1)},\qquad
 J_q=\ker(D^q\longrightarrow C),\qquad
 P^q=\varprojlim_n D^q/J_q^n,
\]
where the map to $C$ uses $B\to C$ and $a$ on every $A$-factor.  The usual coface maps, inserting $1$ in an $A$-factor, define an augmented cosimplicial $B$-algebra
\[
                    B\longrightarrow P^0\rightrightarrows P^1
                              \longrightarrow P^2\longrightarrow\cdots.
\]
For every $B$-module $N$ its alternating augmented complex
\[
 0\longrightarrow N\longrightarrow P^0\otimes_BN
       \longrightarrow P^1\otimes_BN\longrightarrow\cdots
\]
is exact.  In particular the first two arrows exhibit $N$ as the equalizer of the two coface maps from $P^0\otimes_BN$ to $P^1\otimes_BN$.
\end{lemma}
\begin{proof}
Every $D^q$ is a Noetherian ring and is flat over $B$.  Noetherian adic completion is flat, so every $P^q$ is flat over $B$.  Completion is exact on finite modules and commutes with their tensor description; consequently
\[
 P^q/IP^q\simeq
 \widehat{C\otimes_R A^{\otimes_R(q+1)}},
\]
where the completion on the right is along the graph ideal.  These are the precise Noetherian completion facts in \citet[Lemmas~10.97.1--6, especially Tags~00MA and~00MB]{StacksProject}.

Modulo $I$, the augmented formal \v{C}ech nerve has a section.  Explicitly, applying $a$ to the first $A$-factor and multiplying its image into the $C$-factor gives continuous homomorphisms
\[
 C\otimes_R A^{\otimes_R(q+1)}
       \longrightarrow C\otimes_R A^{\otimes_R q}.
\]
They preserve the graph ideals and hence extend to the completions.  For $q=0$ the target is $C$.  These maps are extra codegeneracies: their identities with insertion of $1$ are the usual split augmented \v{C}ech identities.  The resulting alternating complex is therefore contractible after augmentation, and remains so after tensoring over $C$ with any $C$-module.

Filter $N$ by the finite filtration $I^jN$.  Since all $P^q$ are $B$-flat, this induces a filtration by subcomplexes whose successive quotient for $Q_j=I^jN/I^{j+1}N$ is the just-described split complex
\[
 0\longrightarrow Q_j\longrightarrow
 (P^0/IP^0)\otimes_CQ_j\longrightarrow
 (P^1/IP^1)\otimes_CQ_j\longrightarrow\cdots.
\]
Every quotient is exact.  Induction on the finite filtration, using the long exact cohomology sequence, proves exactness of the original complex.  No finiteness condition on $N$ was used.
\end{proof}

\begin{theorem}[Ordinary coherent crystals for flat morphisms over a Noetherian base]\label{thm:DE-E2-coherent-Noetherian}
\status{Proved; root and monodromy full independent proof/source audits passed.}
Let $S$ be locally Noetherian and let $X\to S$ be flat and locally of finite presentation.  Restriction is an equivalence
\[
 \left\{\begin{array}{c}
 \text{quasi-coherent infinitesimal crystals}\\
 \text{with coherent value on }X
 \end{array}\right\}
 \ \simeq\
 \left\{\begin{array}{c}
 \text{ordinary coherent stratified}\\
 \OO_X\text{-modules over }S
 \end{array}\right\}.
\]
The assertion uses ordinary tensor pullback in both categories.  It holds for the small or big infinitesimal site, also when arbitrary classical tests are allowed.  On Noetherian tests the values of these crystals are coherent; on arbitrary tests they are of finite presentation.  No flatness of the coefficient module, and no restriction on the characteristic, is imposed.

Thus the simultaneous locally Noetherian and coherent proviso in the manuscript is valid, although the broader assertion requested in Part~E.2 is false by Theorem~\ref{thm:DE-E2-counterexample} and Corollary~\ref{cor:DE-E2-Noetherian-counterexample}.
\end{theorem}
\begin{proof}
We divide the proof into ordinary full faithfulness and essential surjectivity.  Only the latter uses a derived comparison theorem.  In that step we take zeroth cohomology; we do not assert that the intermediate derived crystal is concentrated in degree zero.

\emph{Full faithfulness on Noetherian tests.}
Let $\mathcal E,\mathcal F$ be the crystals in the statement and let
$f_X:\mathcal E_X\to\mathcal F_X$ be a morphism of their ordinary stratified modules.  Work over $S=\Spec R$ and an affine open $V=\Spec A\subset X$.  On an affine Noetherian test $U=\Spec C\hookrightarrow T=\Spec B$ whose map to $X$ factors through $V$, retain the notation $I,D^q,J_q,P^q$ from the lemma.  In the small site we may take $U=V$ after localization.  Since $I$ is nilpotent, lifting a finite set of generators modulo $I$ shows that both $E_T=\Gamma(T,\mathcal E_T)$ and $F_T=\Gamma(T,\mathcal F_T)$ are finite $B$-modules.  Hence they are finitely presented.

For every $q,n$, the test $U\hookrightarrow\Spec(D^q/J_q^n)$ maps both to $T$ and, using any selected $A$-factor, to $V$.  The crystal identifications, followed by completion of finite modules, give canonical isomorphisms
\begin{equation}\label{eq:DE-E2-formal-test-values}
 P^q\otimes_BE_T\simeq P^q\otimes_AE_V,\qquad
 P^q\otimes_BF_T\simeq P^q\otimes_AF_V.
\end{equation}
Here the chosen $A$-factor specifies the right-hand tensor product.  This passage to completion is justified because $D^q\otimes_BE_T$ and $D^q\otimes_AE_V$, and likewise the two modules for $F$, are finite over the Noetherian ring $D^q$.  No tensor product with an arbitrary inverse limit is interchanged.

Use \eqref{eq:DE-E2-formal-test-values} for $q=0$ to transport $P^0\otimes_Af_V$ to a homomorphism
\[
                 \widetilde f:P^0\otimes_BE_T\longrightarrow P^0\otimes_BF_T.
\]
Its two pullbacks to $P^1$ agree.  Indeed the two choices of $A$-factor in \eqref{eq:DE-E2-formal-test-values} differ by the given stratification: this follows already on each finite graph neighbourhood from the crystal cocycle, and then by completion.  Compatibility of $f_X$ with these stratifications is exactly the required equality.

Put $H=\Hom_B(E_T,F_T)$, a finite $B$-module.  Since $E_T$ is finitely presented and $P^q$ is $B$-flat, the canonical map
\[
 P^q\otimes_BH\ \simeq\
 \Hom_{P^q}(P^q\otimes_BE_T,P^q\otimes_BF_T)
\]
is an isomorphism; it follows by applying $\Hom(-,F_T)$ to a finite presentation of $E_T$ and using flatness to preserve its kernel.  The equalizer assertion of Lemma~\ref{lem:DE-E2-completed-Cech} therefore descends $\widetilde f$ to a unique $B$-linear map $f_T:E_T\to F_T$.

These maps are compatible with morphisms and open restrictions of Noetherian tests.  To check this, pull both proposed maps to the completion $P^0$ for the source test.  The maps of finite graph neighbourhoods, and hence their completions, identify both pullbacks with the given map $f_V$.  The map $B\to P^0$ is faithfully flat: it is flat as above and its composition with $P^0\twoheadrightarrow C$ is the nilpotent quotient $B\twoheadrightarrow C$, whose spectrum is surjective.  The two maps on the test are consequently equal.  The same equalizer uniqueness shows that a morphism of crystals is uniquely determined by its value on $X$.  This proves ordinary full faithfulness on Noetherian tests.

\emph{Lifting a coherent stratification to the derived formal nerve.}
Let $M$ be a coherent ordinary stratified module on $X/S$.  Continue locally with $X=\Spec A$ and $S=\Spec R$.  Let $Q_q$ be the completion of $A^{\otimes_R(q+1)}$ along the intrinsic diagonal.  Each projection $A\to Q_q$ is flat: tensoring the other flat $R$-factors is flat, and Noetherian completion is flat.  Therefore
\[
                    Q_q\otimes_A^{\mathbf L}M=Q_q\otimes_AM[0].
\]
These finite $Q_q$-modules are complete for the diagonal ideal.  The compatible ordinary stratification isomorphisms on the finite diagonal neighbourhoods give isomorphisms between these completed modules, because completion commutes with finite modules.  The identities and cocycle conditions continue to hold after completion.

These are derived stratification data as well.  Namely, derived pullback along a map between terms of the formal nerve is derived completed pullback; on a module of the displayed form it gives the corresponding $Q_r\otimes_A^{\mathbf L}M$, which is again the ordinary module because $Q_r$ is $A$-flat.  All mapping spaces between these modules are zero-truncated, so the ordinary identities and cocycles supply the coherent derived data.  This reasoning concerns the complete coherent modules; it does not assert that the finite diagonal neighbourhoods are flat over $A$.

Jiang's Theorem~2.1 and Remark~2.4 identify this intrinsic derived stratifying category; Proposition~6.2 identifies a formal neighbourhood with derived complete modules (equivalently supported modules); and Theorem~6.3 identifies the category with derived quasi-coherent infinitesimal crystals, compatibly with evaluation and restriction.  The hypotheses hold because $A$ is finitely presented and flat over the Noetherian ring $R$, hence has finite Tor-amplitude.  See \citet[pp.~16--17, 32--34; also Theorem~B.5, pp.~43--45]{Jiang2023Derived}.  We obtain a derived crystal $K$ with $K_X=M[0]$ and with the prescribed completed stratification.

\emph{Zeroth cohomology gives an ordinary coherent crystal.}
For every finite-type test $U=\Spec C\hookrightarrow T=\Spec B$, its derived reduction satisfies
\[
 K_T\otimes_B^{\mathbf L}C=K_U\in D^{\leq0}(C).
\]
Indeed $K_U$ is the derived pullback of the ordinary module $M$ along $U\to X$.  For every $C$-module $Q$, associativity and right t-exactness give
\[
                   K_T\otimes_B^{\mathbf L}Q
                   =K_U\otimes_C^{\mathbf L}Q\in D^{\leq0}(B).
\]
Filter $B$ by the finitely many powers of the nilpotent ideal $I$.  Its successive quotients are $C$-modules.  The exact triangles of this finite filtration show that $K_T\in D^{\leq0}(B)$.  No bounded-below assumption has been used.

Set $G_T=H^0(K_T)$.  If $h:T\to T'$ is a morphism of tests, derived crystal compatibility and right t-exactness imply
\[
 H^0(K_T)=H^0(Lh^*K_{T'})=h^*H^0(K_{T'}).
\]
The last equality follows, for any complex in $D^{\leq0}$, by its truncation triangle: derived pullback sends the $D^{\leq-1}$ term into $D^{\leq-1}$.  Thus the $G_T$ form an ordinary quasi-coherent crystal.  Moreover
\[
          G_T/IG_T=H^0(K_U)=C\otimes_AM,
\]
which is a finite $C$-module.  Nilpotent Nakayama makes $G_T$ finite over $B$, and hence coherent on these Noetherian tests.  On $X$ its value is $M$, and on the finite intrinsic diagonal neighbourhoods its ordinary stratification is exactly the original one: zeroth cohomology of the derived reduction of the completed isomorphism is its ordinary reduction.  This proves essential surjectivity on finite-type tests.

\emph{Arbitrary tests and gluing.}
Every affine test over an affine $V=\Spec A$, where $A=R[z_1,\ldots,z_r]/K_0$, maps to a Noetherian finite-type test
\[
        V\hookrightarrow\Spec\bigl(R[z_1,\ldots,z_r]/K_0^N\bigr)
\]
by lifting the $z_i$ and choosing $N$ large enough to annihilate the nilpotent ideal of the given test.  Any two choices admit a common refinement, obtained from their finite-type product by taking a sufficiently high graph neighbourhood of $V$.  Parallel compatible maps are equalized by their closed equalizer, which is still a Noetherian finite-type thickening of $V$.  Consequently pullback extends the ordinary crystal uniquely, on objects and morphisms, from finite-type tests to arbitrary classical tests.  The values on arbitrary tests are of finite presentation, because they are pullbacks of finite modules over Noetherian rings.

Full faithfulness, already proved locally, makes the local constructions for $M$ glue uniquely on overlaps of affine opens of $X$ and $S$; their cocycle relations follow from the same uniqueness.  The argument works on either infinitesimal site, with the domain of a big test mapping to the chosen affine open.  We have therefore obtained the claimed equivalence with ordinary tensor pullback in the full stated scope.
\end{proof}

\begin{remark}[Why zeroth cohomology is essential]\label{rem:DE-E2-coherent-derived-heart}
\status{Proved; root and monodromy full independent audits passed.}
The proof establishes an equivalence of the two ordinary coherent categories.  It does not identify an ordinary crystal with a derived crystal concentrated in degree zero on every test.  For nonflat coefficients the intermediate $K_T$ can have negative cohomology.  What is proved is the uniform upper bound $K_T\in D^{\leq0}$, which is exactly what permits $H^0$ to commute with ordinary tensor pullback.  The Noetherian and coherent hypotheses enter separately through flatness of completion and completion of finite modules; dropping either leads to the counterexamples already given.
\end{remark}

\subsection{E.3. Integral infinitesimal cohomology without divided powers}
\status{Proved; Ogus's integral comparison, with local embeddings removing the global embedding restriction.}

\begin{theorem}\label{thm:DE-inf}
\status{Proved under stated hypotheses}
Let $X/k$ be smooth proper, with $k$ algebraically closed of characteristic $p$, and $W=W(k)$.  On the infinitesimal site without divided powers, there is a natural integral isomorphism
\[
 H^i((X/W)_{\mathrm{inf}},\OO)
   \cong H^i_{\et}(X,\Zp)\otimes_{\Zp}W
\]
for every $i$.  It includes $p$-power torsion.
\end{theorem}
\begin{proof}
With a global smooth ambient embedding, this is precisely \cite[Corollaries~4.7 and~4.11, Remark~4.10]{Ogus1975}; Theorem~4.9 identifies the image in crystalline cohomology with the Frobenius-stable part.  We spell out why smooth proper schemes without such a global embedding satisfy the same conclusion.

Choose a finite affine open cover $(U_i)$ of $X$ and closed embeddings of the $U_i$ in special fibres of affine spaces $P_i/W$.  Separatedness implies that every multiple intersection $U_I$ embeds as a closed subscheme in $\prod_{i\in I}U_i$, and hence in $P_I=\prod_{i\in I}P_i$.  These embeddings are compatible with the projection maps of the \v Cech nerve.  Combining this affine cover with the \v Cech--Alexander construction in each ambient $P_I$ gives a nonnegative double complex computing the infinitesimal cohomology over $W_n=W/p^n$.  Its terms are finite products of rings of functions on completions of smooth affine $W_n$-schemes along $U_I$.  Write its total complex as $C_n$.

Each such completed ring is flat over $W_n$: it is a noetherian adic completion of a smooth $W_n$-algebra, hence flat over that algebra and over $W_n$.  Adic completion commutes with the indicated finite quotients.  Consequently, term by term, there is an exact sequence
\[
 0\longrightarrow C_1\xrightarrow{p^{n-1}}C_n\longrightarrow C_{n-1}\longrightarrow0.\tag{E.3.1}
\]
The \v Cech--Alexander calculation uses ordinary inverse limits of affine nilpotent neighbourhoods, whose restriction maps are surjective, so no higher inverse limit is lost at this step.

For $n=1$, smoothness makes the relative Frobenius finite flat.  Its iterated equivalence relations are cofinal among the infinitesimal diagonal neighbourhoods.  Faithfully flat descent therefore identifies the characteristic-$p$ \v Cech--Alexander complex with $R\!\varprojlim_F R\Gamma(X,\OO_X)$; this is the argument made explicit in \cite[Remark~4.12]{Ogus1975}.  Since $X$ is proper, every $H^i(X,\OO_X)$ is finite dimensional.  Its Frobenius images stabilize, the inverse systems are Mittag--Leffler, and only the Frobenius-stable part survives.  Artin--Schreier and the semilinear fixed-point lemma over algebraically closed $k$ give
\[
 H^i(C_1)=H^i(X,\OO_X)_{\mathrm{stable}}
     \cong H^i_{\et}(X,\Fp)\otimes_{\Fp}k.
\]
The natural comparison is obtained by pulling the constant coefficient sheaf into the \etale infinitesimal site and then comparing \etale and Zariski cohomology for its coherent structural values, as in \cite[Propositions~1.8 and~1.13]{Ogus1975}.  It is compatible with reduction and multiplication by $p$.

On the \etale side the short exact sequence
$0\to\Fp\to\mathbb Z/p^n\to\mathbb Z/p^{n-1}\to0$, followed by the flat scalar extension to $W$, gives the corresponding exact triangle.  Induction using (E.3.1) and the five lemma therefore gives the finite-level comparison
\[
 H^i(C_n)\cong H^i_{\et}(X,\mathbb Z/p^n)\otimes_{\mathbb Z/p^n}W_n
\]
for all $n$.

These groups are finite-length $W_n$-modules.  Inverse systems of their images in any fixed level stabilize, so the Mittag--Leffler condition eliminates $R^1\!\varprojlim$ on cohomology.  On complexes, the transitions $C_n\to C_{n-1}$ are termwise surjective.  The cohomology of the site over $W$ is consequently the inverse limit of the displayed finite-level groups.  The inverse limit of the finite \etale groups is the usual finitely generated $\Zp$-module $H^i_{\et}(X,\Zp)$, and finite generation makes ordinary tensoring with the complete flat ring $W$ agree with its $p$-adic completion.  Taking the limit proves the theorem.
\end{proof}

\subsection{E.4. The intrinsic completed nerve, including the cusp}
\status{Known theorem; the cusp is computed explicitly.}

\begin{theorem}\label{thm:DE-intrinsic-nerve}
\status{Proved under stated hypotheses}
For every finite type complex scheme $X$, the formal completions of $X^{n+1}$ along the small diagonals compute infinitesimal/de Rham cohomology, when one takes the completed cosimplicial complex and its Zariski hypercohomology.  For the cusp $X=\Spec\mathbb C[t^2,t^3]$, this intrinsic complex has $H^0=\mathbb C$ and $H^i=0$ for every $i>0$.
\end{theorem}
\begin{proof}
The precise primary result is \cite[Corollary~4.30]{Bhatt2012}; the affine assertion is Corollary~4.19 and the comparison with the infinitesimal site is Remark~4.28.  Smoothness of $X$ is not required.  Flatness over $\mathbb C$ is automatic.  In affine notation the statement is
\[
 R\Gamma(X_{\mathrm{inf}},\OO)\simeq
 \operatorname{Tot}\bigl(A\longrightarrow\widehat{A\otimes A}
     \longrightarrow\widehat{A^{\otimes3}}\longrightarrow\cdots\bigr),
\]
where completion is along the multiplication map to $A$ and the differential is the alternating sum of coprojections.  The completed complex must not be replaced by the ordinary Amitsur complex or by the K\"ahler de Rham complex at a singular point.

For the cusp $A=\mathbb C[t^2,t^3]=\mathbb C[x,y]/(y^2-x^3)$, its term of cosimplicial degree $n$ is explicitly
\[
 C^n=\mathbb C[x_0,y_0,\ldots,x_n,y_n]/(y_i^2-x_i^3)_{0\leq i\leq n}
       \ {}^{\widehat{\ }}_{(x_i-x_0,y_i-y_0)_{1\leq i\leq n}}.
\]
The normalization $\mathbb C\to X^{\mathrm{an}}$, $t\mapsto(t^2,t^3)$, is a homeomorphism: it is proper and bijective, and its inverse away from the origin is $y/x$, extending continuously at the origin.  Thus $X^{\mathrm{an}}$ is contractible.  Bhatt's comparison, in its explicit Betti formulation \cite[Example~4.23]{Bhatt2012}, yields
\[
 H^0(C^\bullet)=\mathbb C,\qquad H^i(C^\bullet)=0\quad(i>0).
\]
This is a computation of the intrinsic complex itself, not a replacement of that complex by an unproved singular Poincar\'e lemma.
\end{proof}

\subsection{E.5. K\"ahler de Rham cohomology of affine locally planar curves}
\status{Proved for reduced affine curves, the standard meaning of plane curve singularities.}

\begin{theorem}\label{thm:DE-curves}
\status{Proved under stated hypotheses}
Let $C$ be a reduced affine complex curve with locally planar singularities.  For its ordinary K\"ahler de Rham complex,
\[
 \dim_{\mathbb C}H^1(\Omega_C^\bullet)
       =b_1(C^{\mathrm{an}})+\sum_{x\in C_{\mathrm{sing}}}(\mu_x-\tau_x),
 \qquad H^i(\Omega_C^\bullet)=0\quad(i\geq2).
\]
$H^0$ is the space of locally constant functions, of dimension the number of connected components.
\end{theorem}
\begin{proof}
Bonnet's \cite[Theorem~1.2]{Bonnet2006} gives
\[
 \dim\bigl(\Gamma(C,\Omega_C^1)/d\Gamma(C,\OO_C)\bigr)
       =b_1(C^{\mathrm{an}})+\sum_x\mu_x
\]
for locally complete-intersection curve singularities; his local invariant is the usual Milnor number in this case.  This local identification also follows from \cite[Theorem~2.2(1)]{DimcaGreuel2018}.

At a singular point with completed local equation $f(x,y)=0$, one has
\[
 \widehat{\Omega_C^2}_x\cong
 \frac{\mathbb C[[x,y]]}{(f,f_x,f_y)}\,dx\wedge dy,
\]
so its length is $\tau_x$.  The coherent module $\Omega_C^2$ is supported at the finitely many singular points, and $\Omega_C^3=0$.

We verify global surjectivity of $d:\Gamma\Omega_C^1\to\Gamma\Omega_C^2$; surjectivity as a map of sheaves alone would not suffice, since $d$ is not $\OO_C$-linear.  At every completed local singularity, a two-form represented by $h\,dx\wedge dy$ has a formal primitive $(\int h\,dx)\,dy$, since the characteristic is zero.  Choose $N_x$ with $\mathfrak m_x^{N_x}\Omega_{C,x}^2=0$.  Chinese remainder for the finite module $\Gamma(C,\Omega_C^1)$ allows simultaneous interpolation of the formal primitives modulo $\mathfrak m_x^{N_x+1}\Omega_{C,x}^1$.  An error in that power differentiates into $\mathfrak m_x^{N_x}\Omega_{C,x}^2=0$.  Thus the interpolated global one-form has exactly the prescribed differential at every singularity, hence globally.  It follows that
\[
 0\to H^1(\Omega_C^\bullet)\to
 \Gamma\Omega_C^1/d\Gamma\OO_C\xrightarrow{d}\Gamma\Omega_C^2\to0
\]
is exact and $H^2=0$.  Taking dimensions subtracts $\sum\tau_x$ and gives the formula.
\end{proof}

\subsection{E.6. Integral Monsky--Washnitzer independence}
\status{Known theorem, integrally in all degrees and for every prime.}

\begin{theorem}\label{thm:DE-integral-MW-independence}
\status{Proved under stated hypotheses}
Integral Monsky--Washnitzer cohomology of a smooth finite type algebra over a perfect field of characteristic $p$ is canonically independent of its smooth weak lift and is functorial in the special-fibre algebra.  Two maps between weak lifts that agree modulo $p$ induce the same map in every cohomological degree, without killing torsion.
\end{theorem}
\begin{proof}
The precise integral assertion is known: \cite[Proposition~1]{ErtlSprang2018} proves that if $A^\dagger$ is the weak completion of a smooth $W(k)$-lift of a smooth finite type $k$-algebra and $B$ is any weakly complete $W(k)$-algebra, two maps $A^\dagger\to B$ agreeing modulo $p$ induce chain-homotopic maps of the integral continuous de Rham complexes.  Proposition~2 and Lemma~3 give exactly canonical independence of the smooth lift and functoriality of integral Monsky--Washnitzer cohomology.

The integrality mechanism matters.  The homotopy is constructed in a weak completion with parameter $T$, with endpoints $T=0$ and $T=p$.  Integration of $T^i\,dT$ between these endpoints gives
\[
 \frac{p^{i+1}}{i+1}\in W(k),
\]
whose valuations are nonnegative and tend to infinity rapidly enough for the weak convergence condition.  Thus this is not the generally invalid integral integration between $0$ and $1$.  Very smoothness gives the required interpolating homomorphism.  Given two lifts $A^\dagger,(A')^\dagger$, choose maps in both directions lifting the identity of the special fibre.  Their composites lift the identity and are therefore homotopic to the identity.  The induced cohomological isomorphisms are independent of those choices and mutually inverse.  Maps of special-fibre algebras induce maps in the same direction on these algebraic de Rham complexes, equivalently contravariant maps in schemes.  No torsion is discarded.
\end{proof}

\subsection{E.7. The Lubin--Tate tangent exponent and the two full lattices}
\status{Proved; the parent agent independently audited the complete corrected comparison.  The tangent exponent is zero, while the full lattices can differ.}

Write $A$ for the ring of integers of the local field in the question, $k=\mathbb F_q$, $q=p^f$, and $W=W(k)$.  Let $\pi$ be the chosen uniformizer, and choose an Eisenstein equation $E(u)$ for $\pi$ over $W$, multiplied by a unit so that $E(0)=p$.  Let $\sigma$ act on $\mathfrak S=W[[u]]$ by Witt Frobenius and $u\mapsto u^p$.  In the formulas below, the copy of $A$ acting on the group and the copy of $A$ which is the base must be distinguished.  Thus
\[
 \Lambda=A\otimes_{\mathbb Z_p}\mathfrak S,
 \qquad R=A\otimes_{\mathbb Z_p}A.
\]
The $f$ factors of $A\otimes_{\mathbb Z_p}W$ have idempotents $e_i$, indexed cyclically so that $e_0$ is the identity embedding and $\sigma(e_i)=e_{i-1}$.  Throughout this subsection $\Omega^\vee=\Hom_{\mathbb Z_p}(\Omega,\mathbb Z_p)$, with its induced $A$-module structure; it is not the $A$-linear dual.  In particular $A^\vee$ identifies by the field trace with the inverse different $\mathfrak D_{A/\mathbb Z_p}^{-1}$.

\begin{proposition}\label{prop:DE-LT}
\status{Proved under stated hypotheses}
With the normalization of the crystalline comparison in the question,
\[
                         \alpha(A)=0
\]
for every $A$.  Nevertheless the integral de Rham Dieudonn\'e lattice of the Lubin--Tate group need not equal the scalar extension of the special-fibre Dieudonn\'e lattice.  A counterexample, for every prime $p$, is
\[
                   A=\mathbb Z_p[\pi],\qquad \pi^{p+1}=p.
\]
\end{proposition}
\begin{proof}
We give the integral calculation, including the convention change.  By \cite[Corollary~6.7]{Lau2014}, with the hypotheses verified by his Remark~6.3, Breuil windows over $\mathfrak S\to A$ classify $p$-divisible groups, including when $p=2$.  Take the rank-one $\Lambda$-module $Q=\Lambda$ and the window map
\[
 \phi:Q\longrightarrow Q^{(\sigma)},\qquad
 b=e_0(\pi-u)+(1-e_0).
\]
Here $\pi$ in the first factor is the acting uniformizer.  The quotient is $A$, via $u\mapsto\pi$ and the identity component, and $E(u)$ kills it.  Thus this is a Breuil window.  Its associated group is a one-dimensional formal $A$-module of $A$-height one, with the specified action on its tangent space.  Its special contravariant Frobenius, after one complete cycle of the $f$ residue embeddings, is multiplication by $\pi$.  This is the special Lubin--Tate normalization on the sheet, independently of where the distinguished factor occurs after the covariant Frobenius twist.  A choice of special-fibre identification fixes the lift comparison.  Equivalently, uniqueness can be checked in the window description: the tangent quotient forces the ideal of $\phi$ to be $(b)$, and a factor which is a unit and reduces to $1$ at $u=0$ is removed by the convergent $\sigma$-iteration.  The special-fibre normalization fixes the remaining constant unit.

Lau uses the covariant Dieudonn\'e crystal.  His Proposition~7.1 identifies its integral value on the divided-power envelope of $E$ with the scalar extension of $P=Q^{(\sigma)}$, its Hodge submodule with $\phi(Q)$, and its Frobenius with
\[
                     F_P(x)=\sigma(E/b)\,\sigma(x).
\]
In particular its value at $u=\pi$ is the actual integral de Rham lattice, and its tangent quotient is $P/bQ$.  At $u=0$ its Frobenius is $\sigma(p/b(0))$.  Dualizing gives the contravariant convention of the question: one component has Frobenius factor $\pi$, and the other components have factor $1$, with the cyclic semilinear permutation and the twist in $Q^{(\sigma)}$ understood.  The coordinate-independent statements needed here are that the identity component is the tangent quotient and that the full-cycle contravariant Frobenius is $\pi$.  The scalar factor $\sigma(E)/p$ must be kept in this convention change; omitting it changes the displayed comparison, though not the lattice conclusions below.

The normalized Frobenius-equivariant comparison from the constant special-fibre isocrystal to $P$ is multiplication by
\[
 c(u)=\frac{h(u)}{g(u)},\qquad
 h(u)=\prod_{n\geq1}\frac{\sigma^n E(u)}p,
 \qquad
 g(u)=\prod_{n\geq1}\sigma^n\left(\frac{b(u)}{b(0)}\right).
\]
Indeed $c(0)=1$ and
\[
 c(u)=\frac{\sigma(E)/p}{\sigma(b/b(0))}\,\sigma(c(u)),
\]
which is precisely Frobenius equivariance.  Iteration proves uniqueness.  The products converge on every closed disc of radius less than $1$, after a finite extension containing the coefficients; the exponents $p^n$ give convergence even after division by the fixed uniformizer.  In particular they can be evaluated at the base uniformizer.  The crystalline comparison is the normalized Frobenius-equivariant comparison, so this computation determines it.  One may also check horizontality: the difference from the crystalline connection satisfies the homogeneous Frobenius equation on differentials, whose iterates tend to zero because $d\sigma^n(u)=p^n u^{p^n-1}du$.

Each factor $\sigma^nE(\pi)/p$ is a unit of the base ring $A$.  To see this, the constant coefficient is $1$; every nonleading nonconstant coefficient of $E$ is divisible by $p$, and the valuation of the leading term after division by $p$ is $e(p^n-1)>0$.  Consequently $h(\pi)\in A^\times$.

On the component of the tangent quotient, only the integers $n$ divisible by $f$ contribute to $g$.  Its image under the multiplication map $R\to A$ is therefore
\[
                  \prod_{m\geq1}(1-\pi^{q^m-1})\in A^\times.
\]
The inverse image of the actual lattice $P|_{u=\pi}=R$ in the constant covariant isocrystal is $c(\pi)^{-1}R$.  Its tangent quotient differs from the quotient of the constant lattice by the preceding unit and the unit $h(\pi)^{-1}$.  Thus the two tangent lattices are equal.  To compare with the notation of 8:48--50, choose an $A$-basis of the special covariant line $\Omega^\vee$, thereby identifying it with the rank-one covariant module used above.  Its contravariant dual is then $\Omega$; if the chosen covariant line is $A$, the contravariant line is $A^\vee=\mathfrak D^{-1}$, not canonically $A$.  Replacing this basis tensors both comparison lattices and both tangent quotients by the same invertible $A$-line.  The computed unit ratio is unchanged, so it proves $t_G=\Omega^\vee$ and $\alpha=0$ for the specified special lattice, without discarding the inverse different.

For the full-lattice assertion take $f=1$ and $\pi^{p+1}=p$.  In this case
\[
 g(\pi)=\prod_{n\geq1}
       \left(1-\frac{\pi_{\rm base}^{p^n}}{\pi_{\rm act}}\right)
        \in R[1/p]^\times.
\]
The factors with $n\geq2$ are units of $R$: writing
$\pi_{\rm act}^{-1}=\pi_{\rm act}^{p}/p$, their nonconstant terms are integral and topologically nilpotent because $p^n\geq p^2>p+1$.  The first factor is not in $R$.  Otherwise $\pi_{\rm base}^{p}$ would belong to $\pi_{\rm act}R$, whereas
\[
 R/\pi_{\rm act}R\cong\mathbb F_p[\pi_{\rm base}]/(\pi_{\rm base}^{p+1})
\]
and the image of $\pi_{\rm base}^{p}$ is nonzero.  Hence $g(\pi)\notin R$, while all the remaining factors and $h(\pi)$ are units of $R$.  It follows that $c(\pi)^{-1}R\ne R$.  Dualizing preserves inequality of the two full lattices, proving the assertion for the contravariant crystal $D^*(G)$ as well.
\end{proof}

The lead in \texttt{c7\_c8.md}, H2-3, predicts that the different might force nonzero $\alpha$ at high ramification.  The calculation refutes that prediction.  The full integral comparison does retain ramification information, but the identity tangent quotient discards it.  Thus the earlier reading's phrase ``i.e. whether alpha = 0'' is not an equivalence.  The separate sign correction recorded there is valid: from $t_G=\pi^{-\alpha}\Omega^\vee$ one obtains $\omega_G=\pi^\alpha\Omega$, not the inverse scaling.

\subsection{E.8. Exact transport of the Hasse invariant}
\status{Proved for every abelian scheme over an arbitrary, possibly nonreduced, $\Fp$-scheme; no residual sign.}

Write $H=H^1_{\dR}(A/S)$ and $W_A=f_*\Omega^1_{A/S}$.  The Hodge sequence and the canonical identification $R^1f_*\OO_A=\operatorname{Lie}(A^*)$ give
\[
 0\longrightarrow W_A\longrightarrow H\longrightarrow W_{A^*}^{\vee}
 \longrightarrow0.
\]
Let $\omega_A=\det W_A$.  Fix the determinant convention in which the Hodge subspace is written first: $w_1\wedge\cdots\wedge w_g\wedge q_1\wedge\cdots\wedge q_g$.  Combining the determinant of this sequence with the top de Rham trace defines
\[
 \beta_A:\omega_A\xrightarrow{\sim}\omega_{A^*}.
\]

\begin{theorem}\label{thm:DE-hasse}
\status{Proved under stated hypotheses}
For every prime $p$ and every such abelian scheme,
\[
                  \beta_A^{\otimes(p-1)}(\sigma_A)=\sigma_{A^*}.
\]
\end{theorem}
\begin{proof}
We first give a matrix calculation that identifies the isomorphism precisely, rather than merely invoking the existence of some duality isomorphism.  Work over the universal formal deformation of a fibre over a perfect residue field $k$.  Its base is
$B=W(k)[[t_1,\ldots,t_{g^2}]]$: \cite[Corollary~4.8(i)]{Illusie1985BT} proves this for the $p$-divisible group, and \cite[Theorem~1.2.1]{Katz1981ST} transfers it to formal abelian schemes.  These statements hold for $p=2$ as well.  Only the formal family is needed; no algebraization of an unpolarized universal deformation is asserted.

Choose a lift of Frobenius on $B$, the crystalline module of the reduction modulo $p$, and a basis adapted to the lifted Hodge summand.  Normalize its determinant by the top trace.  Frobenius kills the Hodge summand modulo $p$, and its top exterior power is multiplication by $p^g$ under trace.  Since $B$ is $p$-torsion-free, its linearized matrix has the form
\[
 F=C\begin{pmatrix}pI_g&0\\0&I_g\end{pmatrix},
 \qquad \det C=1.
\]
In particular $C$ is invertible.  Verschiebung consequently has matrix
\[
 V=pF^{-1}=\begin{pmatrix}I_g&0\\0&pI_g\end{pmatrix}C^{-1}.
\]
Let $C_{22}$ be the lower-right $g\times g$ block.  The Hasse invariant of $A$ is the determinant of the upper-left block of $C^{-1}$, reduced modulo $p$.  The Hasse invariant of $A^*$ is the determinant of $C_{22}$, because the dual Hasse map is Frobenius on the Hodge quotient.  Jacobi's complementary-minor identity gives
\[
             \det (C^{-1})_{11}=\frac{\det C_{22}}{\det C}
                             =\det C_{22}.
\]
There is no sign in this complementary principal minor.  The determinant normalization ensures that the basis of the dual Hodge line used in the second calculation is the image under $\beta_A$ of the first.  This proves the asserted equality over $B/p$.

For the stated arbitrary base, descend the abelian scheme and the two sections to a finitely generated $\Fp$-algebra.  At each closed point its residue field is perfect, and the formal deformation above pulls back to the family over the completed local ring.  Thus the difference of the two sections vanishes after every such completion.  Completion is faithfully flat on a noetherian local ring, and a nonzero section of a line bundle has nonzero support at a closed point.  The difference therefore vanishes.  Base change proves the result over the original base, including its nilpotents.
\end{proof}

For comparison, \cite[Lemma~2.1.2, Proposition~2.1.3, Theorem~2.1.4]{Bijakowski2018} proves duality compatibility for the Hasse invariant of an arbitrary $p$-divisible group in characteristic $p$, without a reducedness assumption.  The preceding calculation checks that in the abelian case the power-line isomorphism is exactly the one obtained from the Hodge determinant and trace in 9:22.  A change of determinant convention can multiply $\beta_A$ by a sign, but it does not alter this theorem: that sign has $(p-1)$st power $1$ for odd $p$, and is already $1$ in characteristic $2$.

\subsection{E.9. Biduality and the ordinary determinant comparison}
\status{Proved with explicit Poincar\'e, polarization, cup-product, and determinant conventions.  The literal lead's product rule requires a correction.}

Use the standard positive Poincar\'e biduality $\iota_A:A\xrightarrow{\sim}A^{**}$, and define a polarization associated with an ample line bundle by
\[
                 \lambda_L(x)=t_x^*L\otimes L^{-1}.
\]
Then $\lambda_L^\vee\iota_A=\lambda_L$.  Retain the Hodge-first determinant convention defining $\beta_A$ in E.8.  For an ordinary abelian scheme put
\[
 L_A=T_p(A)_{\et},\qquad L_{A^*}=T_p(A^*)_{\et},\qquad
 \epsilon_g=(-1)^{g(g+1)/2}.
\]
R6, whose independence and integrality arguments have been rechecked, defines
\[
 \alpha_A=\frac{\det T_p(\lambda)_{\et}}{\sqrt{\deg\lambda}}:
                      \det L_A\xrightarrow{\sim}\det L_{A^*}.
\]

\begin{theorem}\label{thm:DE-ordinary-sign}
\status{Proved under stated hypotheses}
For standard Poincar\'e biduality one has exactly
\[
                         \alpha_{A^*}=\alpha_A^{-1}.
\]
Under the ordinary $\Fp$-structures
\[
 \omega_A=(\det L_{A^*}/p)\otimes_{\Fp}\OO_S,
 \qquad
 \omega_{A^*}=(\det L_A/p)\otimes_{\Fp}\OO_S,
\]
the Hodge--trace isomorphism is
\[
                    \beta_A=\epsilon_g\,\overline\alpha_A^{-1}.
\]
Equivalently, the resulting map in the direction requested in 9:29,
$\det(A[p])_{\et}\to\det(A^*[p])_{\et}$, is
$\beta_A^{-1}=\epsilon_g\,\overline\alpha_A$.
\end{theorem}
\begin{proof}
Choose a positive integer $N$ for which the dual polarization
$\mu=N\lambda^{-1}:A^*\to A$ is integral.  Set
$\delta=\sqrt{\deg\lambda}>0$.  Then
$\sqrt{\deg\mu}=N^g/\delta$, whereas
$\det T_p(\mu)_{\et}=N^g(\det T_p(\lambda)_{\et})^{-1}$.
Independence of the polarization in R6 gives the first formula.  No alternating Weil-pairing convention is needed for this argument.  Replacing standard biduality by its negative would multiply this formula by $(-1)^g$; it is that convention change, not an ambiguity in standard biduality, which accounts for the earlier caveat.

We next fix the sign in the second formula.  The ordinary identifications of differentials are obtained from the character pairing and $d\log$.  In the notation of \cite[\S3.3, pp.~161--162]{Katz1981ST}, an element $a^*\in T_p(A^*)_{\et}$ determines the invariant form
$d\log e_A(-,a^*)$ on the multiplicative part of $A$.  For any homomorphism $f:A\to B$, functoriality
\[
                 e_B(fx,y)=e_A(x,f^\vee y)
\]
therefore gives $f^*\omega_B(y)=\omega_A(f^\vee y)$.  Because $\lambda^\vee=\lambda$ under standard biduality, the matrix of $\lambda^*:W_{A^*}\to W_A$ is the matrix
$B=T_p(\lambda)_{\et}$, not its negative.  The fact that
$e_{A^*}(y,x)=e_A(x,y)^{-1}$ is already incorporated into the definition of $\omega_{A^*}$ and must not be counted a second time.

It suffices to calculate over an algebraically closed residue field.  To retain information when $p\mid\delta$, make the calculation over the canonical ordinary lift to $W$.  Existence, compatibility with duals, and lifting of homomorphisms and polarizations are stated in \cite[Theorem~1.4 and the paragraph following it]{DworkOgus1986}; the authors explicitly note which additional characterization needs odd $p$, whereas these existence and functoriality assertions hold for all primes.  On this lift the Hodge summand and unit-root summand split.  Choose Tate bases, hence bases $w_i$ of $W_A$ and $q_j$ of $H^1(A,\OO_A)=\operatorname{Lie}(A^*)$.

Here is the local calculation determining the sign.  If $g_{ij}$ is a transition cocycle of $L$, the differential of
$x\mapsto t_x^*L\otimes L^{-1}$ sends a tangent vector $v$ to the additive Picard cocycle $\iota_v(d\log g_{ij})$.  Thus the coefficient matrix of the Hodge component of $c_1(L)$ in the coherent group $H^1(\Omega^1)$ is $B$, transposed if the opposite matrix convention is used; its determinant is $\det B$.  In the \v Cech de Rham total complex with differential
$D=\delta+(-1)^q d$, where $q$ is \v Cech degree, the product of bidegrees $(a,q)$ and $(a',q')$ has sign $(-1)^{a q'}$.  In particular a global one-form times a class in $H^1(\OO)$ is minus its coherent \v Cech product.  Consequently, in de Rham exterior-algebra notation,
\[
                         c_1(L)=-\sum_{i,j}B_{ij}w_i\wedge q_j.
\]
There are no other slope components on the canonical lift: $c_1(L)$ has Frobenius slope $1$, while $\bigwedge^2$ of the Hodge and unit-root summands has slopes $2$ and $0$, respectively.

Expanding the divided exterior power over $K$ yields
\[
 \frac{c_1(L)^g}{g!}
 =(-1)^g(-1)^{g(g-1)/2}\det(B)\,
      w_1\wedge\cdots\wedge w_g\wedge q_1\wedge\cdots\wedge q_g
 =\epsilon_g\det(B)\,(w_1\wedge\cdots\wedge q_g).
\]
The two signs have different sources: the first is the degree-one \v Cech calculation, and the second moves every $q_i$ past the subsequent $w_j$'s.  Riemann--Roch and the degree formula for a polarization give
\[
 \operatorname{Tr}\bigl(c_1(L)^g/g!\bigr)
                         =\chi(L)=\sqrt{\deg\lambda}=\delta;
\]
see \cite[\S16, Riemann--Roch theorem, p.~150]{Mumford1970}.
If $\beta_A$ has scalar $b$ in the chosen determinant bases, it follows that
$\delta=\epsilon_g\det(B)b$, or
$b=\epsilon_g\delta/\det(B)=\epsilon_g\alpha_A^{-1}$.
This identity is first computed over $K$; both sides are integral units by R6 and the integral Hodge--trace isomorphism, so reduction modulo $p$ is legitimate even when $p\mid\delta$.  No division by $\delta$ in characteristic $p$ is used.

In dimension one this explicitly says $\beta_A=-\overline\alpha_A^{-1}$.  For an elliptic curve the convention above gives
$\lambda_{\OO(0)}(x)=\OO((-x)-(0))$, the negative of the usual Abel--Jacobi map $x\mapsto\OO((x)-(0))$.  This agrees with the degree-one \v Cech sign and fixes the normalization at $g=1$.

Finally, E.8 shows that $\beta_A$ preserves the ordinary $\Fp$-structures.  Both sides of the asserted formula are therefore morphisms of rank-one \etale $\Fp$-sheaves.  Equality at geometric points implies equality of such morphisms over any base, including a nonreduced base.  This proves the theorem in its full relative scope.
\end{proof}

The literal Lead E9 asserts that its scalar $c_g$ is multiplicative in products, so that $c_g=c_1^g$.  With the Hodge-first convention fixed in the question the correct product rule is
\[
                \epsilon_{g+h}=(-1)^{gh}\epsilon_g\epsilon_h.
\]
The product order and the Hodge-first order are respectively
\[
 (W_A,Q_A,W_B,Q_B),\qquad (W_A,W_B,Q_A,Q_B).
\]
Reordering contributes $(-1)^{gh}$. For odd $p$, a product of two ordinary elliptic curves already contradicts the literal rule: $\epsilon_1^2=1$ but $\epsilon_2=-1$. In characteristic $2$ all these signs are $1$. Searches for the normalized determinant map $\alpha_A$ itself, using its polarization determinant and square-root degree formula, did not locate an exact later formulation; the source log records the actual queries. This is a limited search result, not a claim that the construction never appeared elsewhere.

\subsection{E.10. A counterexample on the literal old site}
\status{All three requested comparisons are false on the specified flat site: the affine-line example distinguishes integral torsion, the multiplicative group distinguishes rational affine cohomology, and a supersingular elliptic lift refutes the proper test even modulo torsion.}

Let $\mathcal C$ denote literally the site of 7:18: an object over an affine open $U=\Spec A_0$ of $X_0$ is a flat weakly complete $W$-algebra $B$, together with a surjection $B/pB\to A_0$ having nilpotent kernel, with the stated Zariski coverings.  Its structural sheaf has value $B$.

\begin{proposition}\label{prop:DE-old-site}
\status{Proved under stated hypotheses}
For $X_0=\mathbb A^1_k$,
\[
 H^1(\mathcal C,\OO)[p]=0,
 \qquad
 H^1_{\rm MW}(\mathbb A^1_k/W)[p]\ne0.
\]
Thus test (a) fails integrally, although integral Monsky--Washnitzer cohomology itself is independent of its chosen smooth lift by E.6.
\end{proposition}
\begin{proof}
Flatness gives an exact sequence of sheaves
\[
          0\longrightarrow\OO\xrightarrow p\OO
                    \longrightarrow\OO/p\longrightarrow0.
\]
The object $W\langle x\rangle^\dagger$ covers the final sheaf in the sense that every object admits a morphism to it: choose any lift in $B$ of the coordinate $x$ on its special-fibre open, and use the universal property of weak completion.  Therefore a global section of $\OO/p$ is determined by its value $f(x)\in k[x]$ on this object.

For every $N\geq1$ the algebra
\[
                    B_N=(W[x,t]/(t^N))^\dagger
\]
is a permitted flat object over $\mathbb A^1_k$, via $t\mapsto0$.  The two homomorphisms from the smooth lift given by $x\mapsto x$ and $x\mapsto x+t$ force
\[
                      f(x+t)=f(x)\pmod {t^N}.
\]
As this holds for every $N$, $f$ is constant.  Conversely every constant in $k$ defines such a section, so $H^0(\mathcal C,\OO/p)=k$.  These constants all lift through $W\subset H^0(\mathcal C,\OO)$.  The cohomology exact sequence therefore implies $H^1(\mathcal C,\OO)[p]=0$.

On the integral Monsky--Washnitzer complex, the class $[x^{p-1}dx]$ is killed by $p$ because $d(x^p)=p x^{p-1}dx$, and it is nonzero because an integral primitive would have coefficient $1/p$ at $x^p$.  This proves the inequality.
\end{proof}

There is a precise reason to resist identifying this site with a modern overconvergent site.  In \cite[Definitions~3.1, 3.7; Example~3.8(ii)]{Zhang2019}, enlargements may have arbitrary special-fibre maps to $X_0$, and the cohomological structure sheaf used for rigid cohomology is $B[1/p]$.  The universal enlargements of his Example~4.4 already include
\[
       W\langle x,s\rangle^\dagger/(x^n-ps).
\]
For $X_0$ a point, the kernel of reduction to $k$ contains $s$ and is not nilpotent.  These essential objects are absent from the old site.  Moreover the author explicitly distinguishes $H^i(\OO)\otimes K$ from $H^i(\OO[1/p])$.  Thus his Corollary~5.19 is an exact rational rigid comparison for his analytic sheaf and site; it does not prove either comparison for the literal old site.

\par\smallskip\noindent\textbf{The rational affine comparison also fails.}
The following calculation settles the rational assertion of literal Lead E10; it concerns $H^1(\mathcal C,\OO)\otimes_WK$, with precisely the old objects and coefficients.

\begin{lemma}[Finite weak Nakayama lemma]\label{lem:DE-weak-finite}
\status{Proved under stated hypotheses}
Let $A\to B$ be a homomorphism of weakly complete, finitely generated $W$-algebras.  If $B/pB$ is finite over $A/pA$, then $B$ is finite over $A$.
\end{lemma}
\begin{proof}
Here is a proof, to avoid attributing this finite-module statement to a source that only states the surjectivity case.  Choose finitely many weak generators $y_1,\ldots,y_r$ of $B$ over $A$.  Integrality modulo $p$ supplies a monic polynomial $P_j(T)\in A[T]$ for each $j$, of degree $d_j$, with $P_j(y_j)\in pB$.  Lift the quotient $P_j(y_j)/p$ to a dagger series $G_j$ in the variables $y_1,\ldots,y_r$.  Thus $B$ is a quotient of
\[
 A\langle y_1,\ldots,y_r\rangle^\dagger/
       (P_j(y_j)-pG_j:1\leq j\leq r).
\]
This quotient is generated over $A$ by the finitely many monomials $y^\nu$ with $0\leq\nu_j<d_j$.

For completeness, the convergence needed in that last assertion is as follows.  Present $A$ as a quotient of an integral dagger power-series algebra in finitely many base variables.  Choose overconvergent Banach radii for these variables and for the $y_j$, all strictly greater than $1$ and sufficiently near $1$.  Choose the variable radii sufficiently large relative to the base radius that ordinary successive monic division by the $P_j(y_j)$ is bounded: its remainder operator $R$ has norm at most $1$, and its quotient operators $Q_j$ have norm at most $\rho_j^{-d_j}$.  This is possible because each lower coefficient of $P_j$ is integral and its norm tends to at most $1$ as the base radius tends to $1$.  All $G_j$ converge at the chosen radii.  These radii can be decreased towards $1$ so that
\[
             T(f)=p\sum_jG_jQ_j(f)
             \quad\hbox{has operator norm strictly less than }1.
\]
The division identity is
\[
 f=R(f)+\sum_j(P_j-pG_j)Q_j(f)+T(f).
\]
Iteration gives a convergent Neumann series $\sum_{n\geq0}R(T^nf)$ for the remainder.  Its coefficients converge in the fixed base Banach algebra at radius greater than $1$, hence belong to $A$, and its $y$-degrees remain bounded by the $d_j$.  The corresponding quotient series also converge at these radii.  Thus this is an identity of dagger series, not merely an identity in a $p$-adic completion.  It proves finite generation and hence the lemma.
\end{proof}

\begin{theorem}\label{thm:DE-old-Gm}
\status{Proved under stated hypotheses}
For the literal flat site of 7:18,
\[
 H^1(\mathcal C_{\mathbb G_{m,k}},\OO)=0,
 \qquad
 H^1_{\rm MW}(\mathbb G_{m,k}/W)\otimes_WK
       =K\,[dx/x].
\]
Consequently test (a) fails even after tensoring with $K$, and the rational rigid comparison asserted in literal Lead E10 is false.
\end{theorem}
\begin{proof}
Put $A=W\langle x,x^{-1}\rangle^\dagger$.  Its associated object covers the final sheaf of $\mathcal C$: on every object $B$, a lift of the invertible coordinate $x$ is invertible, since the kernel of $B/pB\to\OO(U)$ is nilpotent and $p$ lies in the Jacobson radical of $B$.  The universal property of weak completion then gives $A\to B$.

An additive $\OO$-torsor is trivial on the object $A$.  Indeed its restriction there is governed by Zariski coverings of the affine weak formal scheme, and \cite[\S2, Theorem~14, p.~10]{Meredith1972} gives $H^1=0$ for its structural sheaf.  Equivalently, one refines the local trivializing coverings of that object by principal weak formal opens and applies that theorem.  After choosing a trivialization, the torsor is therefore represented by a natural additive \v Cech cocycle $u$: on an object $B$ with two maps $a_0,a_1:A\to B$, it has value $u_B(a_0,a_1)\in B$, and
\[
 u_B(a_0,a_2)=u_B(a_0,a_1)+u_B(a_1,a_2).
\]
Changing the trivialization subtracts a coboundary $a_1(f)-a_0(f)$, $f\in A$.

First evaluate on the allowed objects
\[
                P_N=A[t]/(t^N),\qquad
                a_0(x)=x,\quad a_1(x)=x+t.
\]
These give $u(x,t)\in A[[t]]$ with $u(x,0)=0$ and
\[
              u(x,t+s)=u(x,t)+u(x+t,s).
\]
The latter equality follows on the objects with two independent nilpotent variables, and hence in the formal power-series ring.  Differentiating in $s$ at $0$ shows that $\partial_tu(x,t)=a(x+t)$, where $a(x)=\partial_tu(x,0)\in A$.  Write
\[
 a(x)=\sum_{m\in\mathbb Z}a_mx^m,\quad
 c=a_{-1},\quad b_s=a_{s-1}/s\quad(s\ne0).
\]
Over $K$, formal integration gives
\[
 u(x,t)=c\log(1+t/x)+
       \sum_{s\ne0}b_s\bigl((x+t)^s-x^s\bigr).
\]
All coefficients of $u$ are integral.  Its coefficient at $t^nx^{-n}$ is $(-1)^{n-1}c/n$, so $c=0$, by taking $n$ to be arbitrarily large powers of $p$.  For each nonzero integer $s$, the coefficient at $t^nx^{s-n}$ is $b_s\binom{s}{n}$.  Take $n=p^{v_p(s)}$.  The binomial coefficient is a $p$-adic unit, also for negative $s$: in $\mathbb F_p[[z]]$, the first nonconstant coefficient of $(1+z)^s$ occurs at $z^{p^{v_p(s)}}$ and is $s/p^{v_p(s)}$ modulo $p$.  Hence $b_s\in W$.  The coefficients $a_m$ have linear $p$-adic decay in $|m|$, whereas division by $s$ loses at most $\log_p|s|$.  Therefore
\[
                         f(x)=\sum_{s\ne0}b_sx^s\in A.
\]
After subtracting its coboundary, $u$ vanishes on every $P_N$.

It remains to check that these evaluations detect the cocycle on every old object; a formal-diagonal calculation alone would not suffice.  This is local on its underlying open $U$, so take $U$ principal in $\mathbb G_m$, and let $A_U$ be the dagger localization of $A$ at a lift of its denominator.  Both maps $a_0,a_1$ extend to $A_U$.  The preceding vanishing also holds on
$P=A_U[t]/(t^N)$, by localization.  Set $\delta=a_1(x)-a_0(x)\in B$.  Its reduction is nilpotent, so, for every integer $m\geq1$, choose $N$ with $\delta^N\in p^mB$.  Evaluation at $x=a_0(x)$ and $t=\delta$ gives a map $P\to B/p^mB$.  Form the ordinary ring fibre product
\[
                         C=B\times_{B/p^mB}P.
\]
Both $B$ and $P$ are finite over $A_U$: modulo $p$, the first is a finite-type nilpotent thickening of $\OO(U)$ with a chosen section, and the second is finite free; apply Lemma~\ref{lem:DE-weak-finite}.  Since $A_U$ is noetherian, $C\subset B\oplus P$ is a finite $A_U$-module.  It is weakly complete, because evaluations of convergent dagger series preserve the defining equality modulo $p^m$; thus it is a weakly complete algebra of finite type over $W$.  It is $W$-flat as a submodule of $B\oplus P$.

The projection $C\to P$ is surjective with kernel $(p^mB,0)$.  The image of this ideal in $C/pC$ has square zero: the product of two such elements is $p$ times an element of $(p^mB,0)\subset C$, since $2m-1\geq m$.  The kernel of $P/pP\to\OO(U)$ is nilpotent.  Consequently $C/pC\to\OO(U)$ has nilpotent kernel, and $C$ is an allowed old-site object.  The pairs $(a_0,x)$ and $(a_1,x+t)$ define two maps $A\to C$.  Naturality of the cocycle and its vanishing on $P$ imply
\[
                         u_B(a_0,a_1)\in p^mB.
\]
This holds for every $m$, and $B$ is $p$-adically separated.  Hence $u_B(a_0,a_1)=0$.  The torsor is trivial, proving the first assertion.

For the second assertion, termwise integration of an overconvergent Laurent series over $K$ removes every differential term except its residue multiple of $dx/x$; the same linear-decay estimate ensures that the primitive remains overconvergent.  The residue is zero on exact forms.  Thus the displayed one-dimensional Monsky--Washnitzer group is obtained directly.  Its standard identification with rigid cohomology for a smooth affine variety, or the elementary annulus calculation itself, shows that the rational old-site group also differs from rigid cohomology.
\end{proof}

\par\smallskip\noindent\textbf{A smooth proper counterexample.}
\status{Proved and independently audited.}
The failure persists for a smooth proper scheme admitting a smooth proper algebraic lift, and hence persists modulo torsion in test (b).

\begin{proposition}\label{prop:DE-old-proper}
\status{Proved under stated hypotheses}
Let $k=\overline{\mathbb F}_3$, let $W=W(k)$, and let $\mathcal E/W$ be the projective smooth elliptic curve
\[
                         y^2z=x^3-xz^2.
\]
Write $E=\mathcal E_k$.  For the literal old site one has
\[
 H^1(\mathcal C_E,\OO)=0,
 \qquad H^1_{\dR}(\mathcal E/W)\simeq W^2.
\]
In particular test (b) is false, including its qualification ``at least modulo torsion.''
\end{proposition}
\begin{proof}
The discriminant of the displayed Weierstrass equation is $64$, a $3$-adic unit.  The invariant differential on its special fibre is a nonzero scalar multiple of $dx/y$.  In characteristic $3$,
\[
 \frac{dx}{y}=\left(\frac{x}{y}\right)^3dx
                  -\left(\frac1y\right)^3x\,dx.
\]
The Cartier operator kills both $dx$ and $x\,dx$, since it selects powers whose exponent is $2$ modulo $3$.  Thus it kills $H^0(E,\Omega^1)$, and, dually, Frobenius on $H^1(E,\OO_E)$ is zero.  The Artin--Schreier sequence, with surjectivity of $a\mapsto a^3-a$ on the algebraically closed field $k$, gives
$H^1_{\et}(E,\mathbb Z/3\mathbb Z)=0$.  Induction using the coefficient exact sequences gives
\[
 H^1_{\et}(E,\mathbb Z/3^n\mathbb Z)=0\qquad(n\geq1).
\]
The finite-level assertion \cite[Corollary~4.7]{Ogus1975}, whose smooth ambient embedding here is $E\hookrightarrow\mathcal E_{W_n}$, therefore gives
\begin{equation}\label{eq:DE-old-proper-infzero}
                  H^1((E/W_n)_{\rm inf},\OO)=0
                  \qquad(n\geq1).
\end{equation}

We explain carefully why this finite-level vanishing detects torsors on the old weak-formal site.  Let $\mathcal T$ be an additive $\OO$-torsor on $\mathcal C_E$, and restrict it to the weak formal completion $\mathcal E^\dagger$.  This means gluing its restrictions on the affine opens, which are old-site objects; the resulting Zariski $\OO$-torsor corresponds to a rank-two extension
\[
     0\longrightarrow\OO_{\mathcal E^\dagger}
       \longrightarrow V\longrightarrow\OO_{\mathcal E^\dagger}
       \longrightarrow0.
\]
For every $N$, the $N$th algebraic diagonal neighbourhood $P_N$ in $\mathcal E\times_W\mathcal E$ is $W$-flat.  Indeed its successive quotients for the diagonal-ideal filtration are symmetric powers of the locally free $\Omega^1_{\mathcal E/W}$.  The two maps from $P_N^\dagger$ to $\mathcal E^\dagger$ give compatible identifications of the two pulled-back extensions, satisfying the cocycle rule on the corresponding triple neighbourhoods.  Thus $V$ carries an ordinary stratification on all these diagonal neighbourhoods.

Reducing modulo $3^n$ turns weak completion into ordinary reduction and gives a stratified extension $V_n$ on the smooth $W_n$-scheme $\mathcal E_n$.  It defines an infinitesimal crystal on $(E/W_n)_{\rm inf}$.  To spell out this standard smooth-ambient construction: for any nilpotent thickening $U\hookrightarrow T$ in that site, formal smoothness locally supplies a map $T\to\mathcal E_n$ lifting $U\to E$.  Pull back $V_n$.  Two choices form a map into a sufficiently high diagonal neighbourhood, so the stratification identifies the pullbacks; the triple cocycle rule gives effective gluing and independence of choices.  The extension and the associated additive torsor are preserved.  Equation~\eqref{eq:DE-old-proper-infzero} says that this stratified extension is split for every $n$.

In particular the underlying class of $\mathcal T|_{\mathcal E^\dagger}$ vanishes modulo every $3^n$.  Meredith's proper comparison theorem, \cite[\S5, Theorem~A, p.~25; or its projective Theorem~B]{Meredith1972}, identifies
\[
 H^1(\mathcal E^\dagger,\OO)
                  \simeq H^1(\mathcal E,\OO)\simeq W.
\]
This module is $3$-adically separated, so the underlying torsor is trivial.  Fix one global trivialization.  The stratification now has an additive cocycle
\[
                         u_N\in H^0(P_N^\dagger,\OO).
\]
Modulo $3^n$, a horizontal trivialization exists by the preceding infinitesimal vanishing.  Its difference from our fixed trivialization is a global function on $\mathcal E_n$, hence an element of $W_n$.  Such a constant has zero diagonal coboundary.  It follows that $u_N$ is zero modulo $3^n$ for every $n$.  Meredith's same proper comparison makes $H^0(P_N^\dagger,\OO)$ a finite $W$-module; $P_N$ is projective and flat.  Hence it is separated, and $u_N=0$ for every $N$.

Finally these diagonal evaluations detect the cocycle on arbitrary old objects, exactly as in Theorem~\ref{thm:DE-old-Gm}; this final step must be included.  Locally choose a smooth affine lift $A$ of an open of $E$ and two maps $a_0,a_1:A\to B$ inducing the same special-fibre map.  Here and below $A$ denotes its weak completion.  Let $J$ be the diagonal ideal and take
\[
                         P_N=(A\otimes_W^\dagger A)/J^N
\]
where $\otimes_W^\dagger$ denotes the weakly completed tensor product; equivalently, weakly complete the algebraic diagonal neighbourhood of the chosen smooth lift.  This is finite over the first $A$ and $W$-flat.  For each $m$, a sufficiently large $N$ makes the pair $(a_0,a_1)$ factor through $P_N\to B/3^mB$.  Indeed the image of $J$ is nilpotent modulo $3$.  The fibre product
\[
                         B\times_{B/3^mB}P_N
\]
is an allowed old-site object by precisely the finiteness, flatness, and square-zero-kernel argument used above.  It forces the cocycle value in $B$ to lie in $3^mB$ for every $m$, hence to vanish.  All local lifts needed here exist: smooth schemes admit local \etale coordinates, one first lifts through the nilpotent kernel modulo $3$, and Hensel's lemma for weakly complete algebras lifts the resulting simple roots; see \cite[Proposition~2.31 and the preceding very-smooth lifting discussion]{DavisLangerZink2012}.  The same argument on common affine opens proves independence of the chosen smooth chart.  Our trivialization on $\mathcal E^\dagger$ therefore descends to a global section of the old-site torsor.  This proves $H^1(\mathcal C_E,\OO)=0$.

Finally, the Hodge exact sequence of the elliptic scheme has two free rank-one outer terms,
\[
 0\longrightarrow H^0(\mathcal E,\Omega^1_{\mathcal E/W})
  \longrightarrow H^1_{\dR}(\mathcal E/W)
  \longrightarrow H^1(\mathcal E,\OO_{\mathcal E})\longrightarrow0.
\]
Thus its middle term is free of rank two.  This proves the second assertion and the claimed failure after killing torsion.
\end{proof}


\subsection{E.11. The transcendental approach under the stated lifting hypothesis}
\status{Proved; the stronger hypothesis on the complex generic fibre is essential to this argument.}

\begin{theorem}\label{thm:DE-simply-connected-convergent}
\status{Proved under stated hypotheses}
Let $k$ be a perfect field of characteristic $p$ and let $\mathcal X/W(k)$ be proper and smooth, with special fibre $S$.  Suppose a specified embedding $W(k)\hookrightarrow\mathbb C$ makes $\mathcal X_{\mathbb C}$ connected and simply connected.  Put $K=\operatorname{Frac}W(k)$.  Every convergent isocrystal on $S/K$ is constant, and every convergent $F$-isocrystal is the pullback of an $F$-isocrystal over $k$.
\end{theorem}
\begin{proof}
Let $\mathcal E$ be such an isocrystal.  Evaluation on the $p$-adic formal completion of $\mathcal X$ yields a vector bundle with convergent integrable connection on the rigid generic fibre.  The realization is fully faithful; see the equivalences and fully faithful functors in \cite[\S1, (1.13), (1.24)--(1.26)]{EsnaultShiho2018}.  Proper rigid GAGA algebraizes the vector bundle and the connection to a pair $(E,\nabla)$ on $\mathcal X_K$.  After the specified embedding $K\hookrightarrow\mathbb C$, the analytic horizontal sections form a complex local system.  Simple connectivity trivializes it.  Complex proper GAGA then makes the horizontal trivialization algebraic.

This trivialization descends to $K$.  Indeed set
\[
 V=\ker\bigl(H^0(\mathcal X_K,E)\xrightarrow{\nabla}
                H^0(\mathcal X_K,E\otimes\Omega^1_{\mathcal X_K/K})\bigr).
\]
Proper coherent base change and flatness of $\mathbb C/K$ show that $V\otimes_K\mathbb C$ is the full horizontal space after base change.  The evaluation map $V\otimes_K\OO_{\mathcal X_K}\to E$ is an isomorphism after that faithfully flat scalar extension, hence already an isomorphism over $K$, compatible with connection.  Full faithfulness of convergent realization now makes $\mathcal E$ the constant isocrystal associated with $V$.

If $\mathcal E$ carries Frobenius, that structure is a horizontal isomorphism between constant objects.  Since $H^0(\mathcal X_K,\OO)=K$, it is precisely a bijective semilinear map $F:V\to V$; hence the original object is a constant convergent $F$-isocrystal.  An algebraic proper smooth lift and the given field embedding are used as stated.  This does not replace the question by the more general conjecture that trivial \etale fundamental group of the special fibre alone forces constancy: \cite[Introduction]{EsnaultShiho2018} explicitly explains why the existence of a lift does not make the generic fibre simply connected.
\end{proof}

\begin{corollary}[The additional assertion of literal Lead E11]\label{cor:DE-BT-constant}
\status{Proved under stated hypotheses}
Under the same hypotheses, every $p$-divisible group over $S$ is isogenous over $S$ to the pullback of a $p$-divisible group over $k$.
\end{corollary}
\begin{proof}
Let $G/S$ be a $p$-divisible group.  Its rational Dieudonn\'e crystal has an associated convergent $F$-isocrystal; this realization is fully faithful up to isogeny, as recalled and used in \cite[\S5, pp.~50--51, and \S5.3]{DeJong1995}.  The preceding theorem makes it constant, associated with an $F$-isocrystal $(V,F)$ over $k$.

One must still exhibit an integral Dieudonn\'e lattice over $k$; constancy of a rational object alone does not specify one.  Choose a closed point $s\in S$.  Its residue field is a finite separable extension of the perfect field $k$.  Choose a finite Galois extension $k'/k$ containing it.  The fibre $G_s$ supplies a full $W(k')$-lattice
$M'\subset V\otimes_K K'$ stable under $F$ and $V_D=pF^{-1}$, where $K'=\operatorname{Frac}W(k')$.  Intersect its finitely many Galois conjugates:
\[
                  M''=\bigcap_{\gamma\in\operatorname{Gal}(k'/k)}\gamma M'.
\]
This remains a full lattice, is stable under $F$ and $V_D$, and carries the evident semilinear Galois descent datum.  Finite \etale descent for $W(k')/W(k)$ descends it to a finite free $W(k)$-module $M_0\subset V$ with the same operators.  Classical Dieudonn\'e theory over the perfect field $k$, also a special case of \cite[Theorem~4.1.1]{DeJong1995}, produces a $p$-divisible group $G_0/k$ with this Dieudonn\'e module.

Thus the rational Dieudonn\'e crystals of $G$ and $G_0\times_k S$ are isomorphic.  The exact full-faithfulness assertion needed to recover the group morphism is \cite[Corollary~5.1.2, formula~(5.1.2.1)]{DeJong1995}.  It applies because $S$ is quasi-compact and finite type over a perfect field, hence over a field with finite $p$-basis.  It gives an invertible quasi-homomorphism between the groups.  Multiplying by a sufficiently large power of $p$ yields an isogeny.  The erratum \cite{DeJong1998Erratum} replaces the incorrect Lemma~7.1.13.1 and repairs the construction used in the proof; it does not withdraw this corollary.  Both the original theorem and that correction have been opened.
\end{proof}

